\subsection{例}

在本小节中（例 7 和 8 除外），K 表示一个非离散的局部紧交换域；dx 表示 K 的加法群上的一个 Haar 测度。

回忆：\(\mod x = |x|\) 当 \(\mathbf{K} = \mathbf{R}\) 时，\(\mod x = |x|^2\) 当 \(\mathbf{K} = \mathbf{C}\) 时，\(\mod x = |x|_p\) 当 \(\mathbf{K} = \mathbf{Q}_p\) 时。

例 1. --- 一般线性群。

令 A 为代数 \(\mathbf{M}_{n}(\mathbf{K})\)。A 的可逆元群 \(A^{*}\) 正是一般线性群 \(\mathbf{GL}(n,\mathbf{K})\)。对于每个 \(X\in A\)，约化范数 \(\mathrm{Nrd}_{\mathrm{A}/\mathrm{K}}(X)\) 是 det X；因此 \(\mathrm{N}_{\mathrm{A}/\mathrm{K}}(X)=(\det X)^{n}\)（Alg., Ch. VIII, §12, No.~3, Prop. 8；cf.~A, III, §9, No.~3, Example 3）。由于 \(X\mapsto{}^{t}X\) 是 A 到其逆代数上的同构，

\[
\mathrm{N} _ {\mathrm{A} ^ {0} / \mathrm{K}} (X) = \mathrm{N} _ {\mathrm{A} / \mathrm{K}} (^ {t} X) = \det (^ {t} X) ^ {n} = (\det X) ^ {n}.
\]

于是，§1, No.~11 的 Prop. 16 表明，测度

\[
\operatorname{mod} (\det X) ^ {- n} \cdot \bigotimes_ {i, j} d x _ {i j} \quad (X = (x _ {i j}))\tag{3}
\]

是 \(\mathbf{GL}(n,\mathbf{K})\) 上的左、右 Haar 测度。

为确定 \(\mathbf{GL}(n,\mathbf{K})\) 上的相对不变测度，我们将依靠以下引理：

引理 5. --- \(\mathbf{GL}(n,\mathbf{K})\) 到 \(C^{*}\) 的连续表示正是形如 \(X \mapsto \chi(\det X)\) 的映射，其中 \(\chi\) 是 \(K^{*}\) 到 \(C^{*}\) 的连续表示。

这样的映射显然是 \(\mathbf{GL}(n,\mathbf{K})\) 到 \(C^{*}\) 的连续表示。反过来，设 \(\psi\) 是 \(\mathbf{GL}(n,\mathbf{K})\) 到 \(C^{*}\) 的连续表示。对 \(x \in K^{*}\)，令

\[
\widetilde {x} = \left( \begin{array}{c c c c c} x & & & & \\ & 1 & & 0 & \\ & & 1 & & \\ & 0 & & \ddots & \\ & & & & 1 \end{array} \right)
\]

并令 \(\chi(x) = \psi(\widetilde{x})\)。于是，对于每个矩阵 \(X \in \mathbf{GL}(n, K)\)，有 \((\det X^{-1})^{\sim} \cdot X \in \mathbf{SL}(n, K)\)。由于 \(\mathbf{SL}(n, K)\) 是 \(\mathbf{GL}(n, K)\) 的换位子群（A, III, §8, No.~9, Prop. 17 的 Cor.），\(\psi((\det X^{-1})^{\sim} \cdot X) = 1\)，从而

\[
\psi (X) = \psi \big ((\det X) ^ {\sim} \big) = \chi (\det X).
\]

由此，§1, No.~8 的 Prop. 10 的 Cor. 1 表明，\(\mathbf{GL}(n,\mathbf{K})\) 上的相对不变测度除一个常数因子外，都是如下形式的测度

\[
\chi (\det X) \cdot \bigotimes_ {i j} d x _ {i j} \quad (X = (x _ {i j})),\tag{4}
\]

其中 \(\chi\) 是 \(\mathbf{K}^*\) 到 \(\mathbf{C}^*\) 的连续表示。

例 2. --- 仿射群。

对于每个 \(X \in \mathbf{GL}(n, \mathbf{K})\) 和每个 \(x \in K^{n}\)，令 \((X, x)\) 为 \(\xi \mapsto X\xi + x\) 在 \(K^{n}\) 上的仿射线性映射。所有 \((X, x)\) 构成 \(K^{n}\) 的仿射群 G（A, II, §9, No.~4）。平移集 T 是 G 的闭正规子群，典范地同构于 \(K^{n}\)；另一方面，\(\mathbf{GL}(n, \mathbf{K})\) 是 G 的闭子群，而 G 是 \(\mathbf{GL}(n, \mathbf{K})\) 与 \(T = K^{n}\) 的半直积。赋予 G 以这样的（局部紧）拓扑，使得 G 成为 \(\mathbf{GL}(n, \mathbf{K})\) 与 T 的拓扑半直积（GT, III, §2, No.~10）。有

\[
(X, x) = (1, x) \cdot (X, 0).
\]

另一方面，若 \(X \in \mathbf{GL}(n, \mathbf{K})\) 且 \(x \in \mathbf{T}\)，则对于每个 \(\xi \in \mathbf{K}^n\)，

\[
(X, 0) (1, x) (X, 0) ^ {- 1} \xi = X (X ^ {- 1} \xi + x) = \xi + X x = (1, X x) \xi ,
\]

因此，T 的自同构 \((1,x)\mapsto(X,0)(1,x)(X,0)^{-1}\) 的模为 mod(det X)（§1, No.~10, Prop. 15）。根据例 1 以及 §2, No.~9 的备注，测度

\[
\operatorname{mod} (\det X) ^ {- n - 1} \cdot \left(\bigotimes_ {i j} d x _ {i j}\right) \otimes \left(\bigotimes_ {i} d x _ {i}\right) \quad (X = (x _ {i j}), x = (x _ {i}))\tag{5}
\]

是 G 上的左 Haar 测度。另一方面，根据 §2, No.~9 的 Prop. 14，

\[
\Delta_ {\mathrm{G}} ((X, x)) = \Delta_ {\mathbf {G L} (n, K)} (X) \Delta_ {K ^ {n}} (x) (\mathrm{mod} \det X) ^ {- 1},
\]

或

\[
\Delta_ {\mathrm{G}} ((X, x)) = \operatorname{mod} (\det X ^ {- 1}).\tag{6}
\]

因此，G 上的一个右 Haar 测度为

\[
(\mathrm{mod} \det X) ^ {- n} \cdot \left(\bigotimes_ {i j} d x _ {i j}\right) \otimes \left(\bigotimes_ {i} d x _ {i}\right).\tag{7}
\]

例 3. --- 严格三角群。

令 \([1, n]\) 为整数 \(m\) 的集合，其中满足 \(1 \leqslant m \leqslant n\)。令 \(J\) 为 \([1, n] \times [1, n]\) 的一个子集，满足以下条件：

\begin{enumerate}
\def\labelenumi{\arabic{enumi})}
\item
  若 \((i,j)\in \mathbf{J}\)，则 \(i <   j\)
\item
  若 \((i,j)\notin J\)，则对每个整数 \(k\) 满足 \(i < k < j\)，\((i,k)\) 和 \((k,j)\) 这两个有序对中至少有一个不属于 \(J\)。
\end{enumerate}

令 \(T_J\) 为如下矩阵的集合：\(Z = (z_{ij})_{1 \leqslant i \leqslant n, 1 \leqslant j \leqslant n}\) 的元素属于 \(K\)，且 \(z_{ii} = 1\)，并且 \(z_{ij} = 0\)，当 \(i \neq j\) 且 \((i,j) \notin J\) 时。这是 \(\mathbf{GL}(n,K)\) 的闭子集。映射 \(Z \mapsto (z_{ij})_{(i,j) \in J}\) 是从 \(T_J\) 到 \(K^s\) 的同胚（其中 \(s\) 表示 \(J\) 的元素个数）。若 \(Z' = (z_{ij}') \in T_J\)，则 \(Z'Z = (z_{ij}'')\)，其中

\[
\begin{array}{c} z _ {i j} ^ {\prime \prime} = z _ {i j} + z _ {i j} ^ {\prime} + \sum_ {i <   h <   j} z _ {i h} ^ {\prime} z _ {h j} \quad \text {for} i <   j, \\ z _ {i j} ^ {\prime \prime} = 0 \text {for} i > j, \quad z _ {i i} ^ {\prime \prime} = 1, \end{array}
\]

从而 \(Z^{\prime}Z \in T_{J}\)。若将 \(T_{J}\) 与 \(K^{s}\) 认同，则映射 \(Z \mapsto Z^{\prime}Z\)（\(Z^{\prime}\) 固定时）对应于一个仿射映射，其行列式为 1；将 \((i,j) \in J\) 中的有序对按字典序排列并应用以下引理即可看出这一点：

引理 6. --- 令 L 为一个全序有限集。对每个 \(\lambda \in L\)，令 \(V_{\lambda}\) 为交换环 \(k\) 上的有限维自由模；对于 L 中的 \(\lambda, \mu\)，若满足 \(\lambda \leqslant \mu\)，令 \(f_{\lambda \mu} \in \operatorname{Hom}_k(V_\mu, V_\lambda)\)。则线性映射

\[
(v _ {\lambda}) _ {\lambda \in \mathrm{L}} \mapsto \left(\sum_ {\mu \geqslant \lambda} f _ {\lambda \mu} (v _ {\mu})\right) _ {\lambda \in \mathrm{L}},
\]

从 \(\prod_{\lambda \in L}V_{\lambda}\) 到 \(\prod_{\lambda \in L}V_{\lambda}\)，其行列式为 \(\prod_{\lambda \in L}\det f_{\lambda \lambda}\)。

立即可化归到 L 是整数区间的情形，此时引理由 A, III, §8, No.~6 的公式 (31) 得出。

若 \(Z \in \mathrm{T}_{\mathrm{J}}\)，则可见存在 \(Z' \in \mathrm{T}_{\mathrm{J}}\) 使得 \(Z'Z = I_n\)，从而 \(Z' = Z^{-1}\)。因此，\(\mathrm{T}_{\mathrm{J}}\) 是 \(\mathbf{GL}(n,\mathrm{K})\) 的闭子群。另一方面，§1, No.~10 的 Prop. 15 表明，测度

\[
\bigotimes_ {(i, j) \in J} d z _ {i j}
\]

是 \(T_{J}\) 上的左 Haar 测度。通过计算 \(ZZ'\)，同理可见此测度也是 \(T_{J}\) 上的右 Haar 测度。

如果在 \(T_{J}\) 的定义中交换行与列的作用，则有类似的结果。

当 J 是有序对 \((i,j)\) 的集合，且满足 i\textless j 时，群 \(T_{J}\) 称为 K 上 n 阶上严格三角群，记为 \(\mathrm{T}_{1}(n,\mathrm{K})\)。其转置称为下严格三角群。

例 4. --- 大三角群。

令 \(n_1, \ldots, n_r\) 为满足 \(\geqslant 1\) 的整数。置 \(p_k = n_1 + \ldots + n_{k-1}\)，\(n = p_{r+1} = n_1 + \cdots + n_r\)。令 \(I_k\) 为整数 \(j\) 的集合，其中满足 \(p_k < j \leqslant p_{k+1}\)，并令 \(J\) 为所有 \(I_k \times I_l\)（其中 \(k < l\)）的并。令 \(G\) 为 \(\mathbf{GL}(n, K)\) 的闭子群，其元素是满足下列条件的矩阵 \((Z_{kl})_{1 \leqslant k \leqslant r, 1 \leqslant l \leqslant r}\)：

\begin{enumerate}
\def\labelenumi{\arabic{enumi})}
\item
  每个 \(Z_{kl}\) 是矩阵 \((z_{ij})_{i\in \mathrm{I}_k,j\in \mathrm{I}_l}\)，其元素属于 \(\mathbf{K}\)，有 \(n_k\) 行和 \(n_l\) 列；
\item
  \(Z_{kl} = 0\)，当 \(k > l\) 时；
\item
  \(Z_{kk} \in \mathbf{GL}(n_k, \mathrm{K})\)，其中 \(1 \leqslant k \leqslant r\)。
\end{enumerate}

分块乘法公式

\[
\begin{array}{r} \left( \begin{array}{c c c c} Z _ {1 1} & 0 & \ldots & 0 \\ 0 & Z _ {2 2} & \ldots & 0 \\ \vdots & \vdots & \ddots & \vdots \\ 0 & 0 & \ldots & Z _ {r r} \end{array} \right) \left( \begin{array}{c c c c} 1 & Z _ {1 2} & \ldots & Z _ {1 r} \\ 0 & 1 & \ldots & Z _ {2 r} \\ \vdots & \vdots & \ddots & \vdots \\ 0 & 0 & \ldots & 1 \end{array} \right) \\ = \left( \begin{array}{c c c c} Z _ {1 1} & Z _ {1 1} Z _ {1 2} & \ldots & Z _ {1 1} Z _ {1 r} \\ 0 & Z _ {2 2} & \ldots & Z _ {2 2} Z _ {2 r} \\ \vdots & \vdots & \ddots & \vdots \\ 0 & 0 & \ldots & Z _ {r r} \end{array} \right) \end{array}\tag{8}
\]

表明 G 是由元素 \((Z_{kl}) \in G\)（满足 \(Z_{kl} = 0\)，其中 \(k \neq l\)）所成的子群 D 与例 3 的子群 \(T_{J}\) 的拓扑半直积。此外，D 同构于各群 \(\mathbf{GL}(n_{k}, \mathbf{K})\)（\(1 \leqslant k \leqslant r\)）的直积。

令 \(J'\) 为有序对 \((j, i)\) 的集合，其中 \((i, j) \in J\)，并令 H 为 \((i, j) \in [1, n] \times [1, n]\) 中不属于 \(J'\) 的有序对的集合。令 \(Z' = (z_{ij})_{1 \leqslant i \leqslant n, 1 \leqslant j \leqslant n}\) 为 G 的一个元素。根据 §2, No.~9 的 Prop. 14 以及上面的例 1 和例 3，取以下测度在映射下的像即可得到 G 上的一个左 Haar 测度：

\[
\bigotimes_ {k = 1} ^ {r} \left(\left(\operatorname{mod} \det Z _ {k k}\right) ^ {- n _ {k}} \cdot \bigotimes_ {i, j \in I _ {k}} d z _ {i j}\right) \otimes \left(\bigotimes_ {(i, j) \in J} d z _ {i j}\right)
\]

在映射

\[
\bigl ((Z _ {k k}), (Z _ {k l}) \bigr) \mapsto \left( \begin{array}{c c c c} Z _ {1 1} & Z _ {1 1} Z _ {1 2} & \ldots & Z _ {1 1} Z _ {1 r} \\ 0 & Z _ {2 2} & \ldots & Z _ {2 2} Z _ {2 r} \\ \vdots & \vdots & \ddots & \vdots \\ 0 & 0 & \ldots & Z _ {r r} \end{array} \right).
\]

现在，对 \(k < l\)，考虑矩阵 \(Z_{kl} = (z_{ij})_{i \in \mathbf{I}_k, j \in \mathbf{I}_l}\) 所成的向量空间。它是 \(n_l\) 个子空间 \(\mathbf{M}_j\)（\(j \in \mathbf{I}_l\)）的直和；这些子空间由满足 \(z_{ih} = 0\)（\(h \neq j\)）的矩阵组成。每个子空间 \(\mathbf{M}_j\) 在映射 \(Z_{kl} \mapsto Z_{kk}Z_{kl}\) 下稳定，并且该映射在 \(\mathbf{M}_j\) 上的限制具有矩阵 \(Z_{kk}\)。因此（§1, No.~10, Prop. 15），测度

\(\bigotimes_{i \in \mathbf{I}_k, j \in \mathbf{I}_l} dz_{ij}\) 在映射 \(Z_{kl} \mapsto Z_{kk} Z_{kl}\) 下的像为

\[
(\mathrm{mod} \det Z _ {k k}) ^ {- n _ {l}} \cdot \bigotimes_ {i \in \mathrm{I} _ {k}, j \in \mathrm{I} _ {l}} d z _ {i j}.
\]

因此，G 上的一个左 Haar 测度为

\[
\prod_ {k = 1} ^ {r} (\mathrm{mod} \det Z _ {k k}) ^ {- q _ {k}} \cdot \bigotimes_ {(i, j) \in \mathrm{H}} d z _ {i j}\tag{9}
\]

其中 \(q_{k} = \sum_{k\leqslant l\leqslant r}n_{l} = n - p_{k}\)

我们再次利用 §2 的 Prop. 14 来计算 G 的模。群 D 和 T \(_{J}\) 是幺模的；另一方面：

\[
\begin{array}{c} \left( \begin{array}{c c c c} Z _ {1 1} & 0 & \ldots & 0 \\ 0 & Z _ {2 2} & \ldots & 0 \\ \vdots & \vdots & \ddots & \vdots \\ 0 & 0 & \ldots & Z _ {r r} \end{array} \right) \left( \begin{array}{c c c c} 1 & Z _ {1 2} & \ldots & Z _ {1 r} \\ 0 & 1 & \ldots & Z _ {2 r} \\ \vdots & \vdots & \ddots & \vdots \\ 0 & 0 & \ldots & 1 \end{array} \right) \left( \begin{array}{c c c c} Z _ {1 1} & 0 & \ldots & 0 \\ 0 & Z _ {2 2} & \ldots & 0 \\ \vdots & \vdots & \ddots & \vdots \\ 0 & 0 & \ldots & Z _ {r r} \end{array} \right) ^ {- 1} \\ = \left( \begin{array}{c c c c} 1 & Z _ {1 2} ^ {\prime} & \ldots & Z _ {1 r} ^ {\prime} \\ 0 & 1 & \ldots & Z _ {2 r} ^ {\prime} \\ \vdots & \vdots & \ddots & \vdots \\ 0 & 0 & \ldots & 1 \end{array} \right), \end{array}
\]

其中 \(Z_{kl}^{\prime}=Z_{kk}Z_{kl}Z_{ll}^{-1}\)。考虑例 3 以及 §1, No.~10 的 Prop. 15，并像上面那样论证可见，若 \(X=\mathrm{diag}(Z_{11},\ldots,Z_{rr})\in\mathbf{D}\)，则自同构 \(Z\mapsto X^{-1}ZX\) 的模（在 \(T_{J}\) 上）为

\[
\prod_ {k <   l} (\mathrm{mod} \det Z _ {k k}) ^ {- n _ {l}} (\mathrm{mod} \det Z _ {l l}) ^ {n _ {k}},
\]

因此

\[
\Delta_ {\mathrm{G}} (Z) = \prod_ {k = 1} ^ {r} (\mathrm{mod} \det Z _ {k k}) ^ {n + n _ {k} - 2 q _ {k}}.\tag{10}
\]

G 的转置群 \(G'\) 以同样方式研究。对于 \(G'\)，可得左 Haar 测度

\[
\prod_ {k = 1} ^ {r} (\mathrm{mod} \det Z _ {k k}) ^ {- p _ {k + 1}} \cdot \bigotimes_ {(j, i) \in \mathrm{H}} d z _ {i j},
\]

以及模

\[
\prod_ {k = 1} ^ {r} (\mathrm{mod} \det Z _ {k k}) ^ {n + n _ {k} - 2 p _ {k + 1}}.
\]

特别地，若取 \(n_{1}=\ldots=n_{r}=1\)，则所得的群 G 是 \(T(n,K)^{*}\)——即子代数 \(\mathbf{M}_{n}(\mathbf{K})\) 的可逆元群，该子代数由矩阵 \(X=(x_{ij})\) 构成，并满足 \(x_{ij}=0\) 当 i\textgreater j 时。这个代数记为 \(\mathrm{T}(n,\mathrm{K})\)，称为上三角代数，而群 \(\mathrm{T}(n,\mathrm{K})^{*}\) 称为 K 上 n 阶上大三角群。于是前面的公式成为：\(\mathrm{T}(n,\mathrm{K})^{*}\) 上的一个左 Haar 测度为

\[
\prod_ {i = 1} ^ {n} (\mathrm{mod} z _ {i i}) ^ {i - n - 1} \cdot \bigotimes_ {i \leqslant j} d z _ {i j} \quad (Z = (z _ {i j}))\tag{9 bis}
\]

而 \(\mathrm{T}(n,\mathbf{K})^*\) 的模为

(10 bis)

\[
\Delta_ {\mathrm{T} (n, \mathrm{K}) ^ {*}} (Z) = \prod_ {i = 1} ^ {n} (\mathrm{mod} z _ {i i}) ^ {2 i - n - 1} \quad \left(Z = \left(z _ {i j}\right)\right).
\]

对于 \(\mathrm{T}(n,\mathbf{K})^*\) 的转置群，即下大三角群，可得左 Haar 测度

\[
\prod_ {i = 1} ^ {n} (\mathrm{mod} z _ {i i}) ^ {- i} \cdot \bigotimes_ {i \geqslant j} d z _ {i j},
\]

以及模

\[
\prod_ {i = 1} ^ {n} (\mathrm{mod} z _ {i i}) ^ {n + 1 - 2 i}.
\]

备注。--- 群 \(\mathrm{T}(n,\mathrm{K})^{*}\) 是 \(\mathbf{GL}(n,\mathrm{K})\) 的闭子群，且 \(\Delta_{\mathrm{T}(n,\mathrm{K})^{*}}\left((z_{ij})\right)=\prod_{i=1}^{n}(\mathrm{mod}z_{ii})^{2i-n-1}\)。我们在例 1 中看到 \(\Delta_{\mathbf{GL}(n,\mathrm{K})}=1\)。若 n\textgreater1，函数

\[
\Delta_ {\mathrm{T} (n, \mathrm{K}) ^ {*}} / \Delta_ {\mathbf {G L} (n, \mathrm{K})}
\]

在 \(\mathrm{T}(n,\mathbf{K})^*\) 上不能延拓为 \(\mathbf{GL}(n,\mathbf{K})\) 到 \(\mathbf{C}^*\) 的连续表示（因为根据引理 5，这样的表示在 \(\mathbf{SL}(n,\mathbf{K})\) 上应等于 1，而 \(\mathrm{mod}(z_{11})^{1 - n}\neq 1\) 对适当选取的 \(z_{11}\) 成立）。由此可知，齐性空间 \(\mathbf{GL}(n,\mathbf{K}) / \mathrm{T}(n,\mathbf{K})^*\) 当 \(n > 1\) 时不存在相对不变测度（\(\S 2\)，No.~6, Th. 3 的 Cor. 1）。

当 n = 2 时，这个齐性空间可认同为 K 上的射影直线。事实上，令 \((e_{1}, e_{2})\) 为 \(K^{2}\) 的典范基。群 \(\mathbf{GL}(2, \mathbf{K})\) 在去掉 0 的 \(K^{2}\) 的直线集上传递地作用，而 \(Ke_{1} - \{0\}\) 的稳定子群是 \(T(2, K)^{*}\)。

例 5. --- 特殊三角群。

我们重新采用例 4 开头的记号，并考虑子群 \(\mathbf{G}_{1} = \mathbf{G} \cap \mathbf{SL}(n, \mathbf{K})\)。这个子群是群 \(\mathbf{D}_{1} = \mathbf{D} \cap \mathbf{SL}(n, \mathbf{K})\) 与 \(T_{J}\) 的拓扑半直积。群 \(D_{1}\) 有一个同构于 \(\mathbf{SL}(n_{r}, \mathbf{K})\) 的正规子群 A，即由元素 \(\text{diag}(Z_{kk})\)（其中 \(Z_{kk} = 1\) 对 k \textless{} r）组成的子群。下面的同态

\[
\varphi : \operatorname{diag} (Z _ {1 1}, \dots , Z _ {r r}) \mapsto (Z _ {1 1}, \dots , Z _ {r - 1, r - 1})
\]

从 \(D_{1}\) 到 \(\mathbf{GL}(n_{1},\mathbf{K})\times\cdots\times\mathbf{GL}(n_{r-1},\mathbf{K})\) 是满射且核为 A。另一方面，\(\varphi\) 是连续的。考虑附录 I 的引理 2，可将 \(D_{1}/A\) 与 \(\mathbf{GL}(n_{1},\mathbf{K})\times\cdots\times\mathbf{GL}(n_{r-1},\mathbf{K})\) 认同。我们将 A 的 Haar 测度记为 \(\mu\)（cf.~例 6），并将

\[
\alpha = \bigotimes_ {k = 1} ^ {r - 1} \left(\left(\operatorname{mod} \det Z _ {k k}\right) ^ {- n _ {k}} \cdot \bigotimes_ {i, j \in I _ {k}} d z _ {i j}\right) \otimes^ {\prime} d \mu \left(Z _ {r r}\right)
\]

满足以下条件的 \(\mathbf{D}_1\) 上的 Haar 测度

\[
\alpha / \mu = \bigotimes_ {k = 1} ^ {r - 1} \left(\left(\operatorname{mod} \det Z _ {k k}\right) ^ {- n _ {k}} \cdot \bigotimes_ {i, j \in I _ {k}} d z _ {i j}\right)
\]

（§2, No.~7, Prop. 10）。然后如例 4 所示，可得 \(\mathbf{G}_1\) 上的一个左 Haar 测度为

\[
\begin{array}{l} \bmod \Big (\prod_ {k = 1} ^ {r - 1} (\det Z _ {k k}) ^ {n _ {k} - q _ {k}} \Big) \\ \qquad \cdot \left[ \bigotimes_ {k = 1} ^ {r - 1} \big ((\bmod \det Z _ {k k}) ^ {- n _ {k}} \cdot \bigotimes_ {i, j \in I _ {k}} d z _ {i j} \big) \otimes^ {\prime} d \mu (Z _ {r r}) \right] \otimes \bigotimes_ {(i, j) \in J} d z _ {i j}. \end{array}
\]

由于 \(\mathbf{G}_1\) 是 \(\mathbf{G}\) 的正规子群，\(\mathbf{G}_1\) 的模是 \(\mathbf{G}\) 的模的限制（§2, No.~7, Prop. 10 b)）。

若 \(n_{r}=1\)，子群 A 退化为单位元，G 上的一个左 Haar 测度为

\[
\operatorname{mod} \left(\prod_ {k = 1} ^ {r - 1} (\det Z _ {k k}) ^ {- q _ {k}}\right) \cdot \bigotimes_ {k = 1} ^ {r - 1} \left(\bigotimes_ {i, j \in I _ {k}} d z _ {i j}\right) \otimes \bigotimes_ {(i, j) \in J} d z _ {i j}.
\]

若取 \(n_{1}=n_{2}=\ldots=n_{r}=1\)，所得群 \(G_{1}\) 称为上特殊三角群，其转置 \(G_{1}^{\prime}\) 称为下特殊三角群。\(G_{1}\) 上的一个左 Haar 测度为

\[
\operatorname{mod} \left(\prod_ {i = 1} ^ {n - 1} z _ {i i} ^ {i - n - 1}\right) \cdot \left(\bigotimes_ {i = 1} ^ {n - 1} d z _ {i i}\right) \otimes \left(\bigotimes_ {1 \leqslant i <   j \leqslant n} d z _ {i j}\right)\tag{11}
\]

而 \(\mathbf{G}_1\) 的模为

\[
\operatorname{mod} \left(\prod_ {i = 1} ^ {n - 1} z _ {i i} ^ {2 i - 2 n}\right).\tag{12}
\]

对于 \(G_{1}^{\prime}\)，同样可得左 Haar 测度

\[
\operatorname{mod} \left(\prod_ {i = 1} ^ {n - 1} z _ {i i} ^ {n - i - 1}\right) \cdot \left(\bigotimes_ {i = 1} ^ {n - 1} d z _ {i i}\right) \otimes \left(\bigotimes_ {1 \leqslant j <   i \leqslant n} d z _ {i j}\right)
\]

以及模

\[
\operatorname{mod} \biggl (\prod_ {i = 1} ^ {n - 1} z _ {i i} ^ {2 n - 2 i} \biggr).
\]

例 6. --- 特殊线性群。

\(\mathrm{T}_{1}(n,\mathrm{K})\) 和 \({}^{t}(\mathrm{T}(n,\mathrm{K})^{*})\) 是 \(\mathbf{GL}(n,\mathbf{K})\) 的闭子群，且二者的交为 \(\{e\}\)。因此，映射 \((M,N)\mapsto M\cdot N\) 是一个连续双射 \(\varphi\)，其定义域为 \(\mathrm{T}_{1}(n,\mathrm{K})\times{}^{t}(\mathrm{T}(n,\mathrm{K})^{*})\)，值域为子集 \(\Omega\) 的 \(\mathbf{GL}(n,\mathbf{K})\)。

引理 7. --- a) 令 \(U = (u_{ij}) \in \mathbf{GL}(n, K)\)。要使 \(U \in \Omega\)，充分必要条件是 \(\det(u_{ij})_{k \leqslant i,j \leqslant n} \neq 0\) 对 \(k = 2, 3, \ldots, n\) 成立。

\begin{enumerate}
\def\labelenumi{\alph{enumi})}
\setcounter{enumi}{1}
\item
  \(\Omega\) 是 \(\mathbf{GL}(n,\mathbf{K})\) 的开子集。
\item
  映射 \(\varphi\) 是从 \(\mathrm{T}_1(n,\mathrm{K})\times {}^t (\mathrm{T}(n,\mathrm{K})^*)\) 到 \(\Omega\) 的同胚。
\end{enumerate}

要使 \(U \in \Omega\)，充分必要条件是存在 \(Z = (z_{ij}) \in \mathrm{T}_1(n,\mathrm{K})\)，使得 \(ZU \in {}^t (\mathbf{T}(n,\mathbf{K}))\)（此时必然

\(ZU \in {}^{t}(\mathrm{T}(n,\mathrm{K})^{*})\)，因为 \(U\) 和 \(Z\) 都可逆）。由前面的结论，若 \(Z\) 存在，则它是唯一的。因此，要使 \(U \in \Omega\)，充分必要条件是线性方程组

\[
\sum_ {k = 1} ^ {n} z _ {i k} u _ {k j} = 0 \quad (1 \leqslant i <   j \leqslant n)
\]

（其中 \((z_{ij})\in \mathrm{T}_1(n,\mathbf{K})\)）有唯一解。这个方程组可以写成

\[
\sum_ {k = i + 1} ^ {n} z _ {i k} u _ {k j} = - u _ {i j} \quad (1 \leqslant i <   j \leqslant n).\tag{13}
\]

固定 \(i\)，关于 \(n - i\) 个未知数 \(z_{i,i+1}\)、\(z_{i,i+2}\)、\ldots、\(z_{i,n}\) 有一个方程组；这些方程组有唯一解的充分必要条件是

\[
\det (u _ {k j}) _ {i + 1 \leqslant k \leqslant n, i + 1 \leqslant j \leqslant n} \neq 0
\]

对 \(i = 1,2,\ldots ,n - 1\) 成立。这证明了 a)。由此可知 \(\Omega\) 在 \(\mathbf{GL}(n,\mathbf{K})\) 中是开集。另一方面，利用克拉默公式求解方程组 (13) 时，\(z_{ij}\) 是 \(u_{ij}\) 的有理函数，且在 \(\Omega\) 中分母非零，因此 \(Z\) 连续依赖于 \(U\)（在 \(\Omega\) 中），这证明了 c)。

现在令 \(\mathbf{G}_{1}^{\prime}\subset{}^{t}(\mathrm{T}(n,\mathrm{K})^{*})\) 为下特殊三角群。映射 \((M,N)\mapsto M\cdot N\) 是连续双射 \(\psi\)，其定义域为 \(\mathrm{T}_{1}(n,\mathrm{K})\times\mathbf{G}_{1}^{\prime}\)，值域为子集 \(\Omega^{\prime}\) 的 \(\mathbf{SL}(n,\mathbf{K})\)。

引理 8. --- a) 令 \(U = (u_{ij}) \in \mathbf{SL}(n, \mathbf{K})\)。要使 \(U \in \Omega'\)，充分必要条件是 \(\det(u_{ij})_{k \leqslant i,j \leqslant n} \neq 0\) 对 \(k = 2,3,\ldots,n\) 成立。

\begin{enumerate}
\def\labelenumi{\alph{enumi})}
\setcounter{enumi}{1}
\item
  \(\Omega'\) 是 \(\mathbf{SL}(n,\mathbf{K})\) 的开子集。
\item
  映射 \(\psi\) 是从 \(\mathrm{T}_1(n,\mathbf{K})\times \mathrm{G}_1'\) 到 \(\Omega^{\prime}\) 的同胚。事实上，令 \(M\in \mathrm{T}_1(n,\mathbf{K})\) 且 \(N\in {}^t (\mathrm{T}(n,\mathbf{K})^*)\)。要使 \(M\cdot N\in \mathbf{SL}(n,\mathbf{K})\)，充分必要条件是 \(N\in \mathbf{G}_1'\)。因此 \(\Omega^{\prime} = \mathbf{SL}(n,\mathbf{K})\cap \Omega\)，引理 8 立即由引理 7 得出。
\end{enumerate}

命题 6. --- a) 群 \(\mathbf{SL}(n,\mathbf{K})\) 是幺模的。

\begin{enumerate}
\def\labelenumi{\alph{enumi})}
\setcounter{enumi}{1}
\item
  令 \(\mu_1\) 和 \(\mu_2\) 分别为上严格三角群 \(\mathrm{T}_1(n, \mathrm{K})\) 和下特殊三角群 \(\mathrm{G}_1'\) 上的左 Haar 测度。\(\mu_1 \otimes \mu_2\) 在同胚 \((M, N) \mapsto M \cdot N^{-1}\) 下的像（该同胚从 \(\mathrm{T}_1(n, \mathrm{K}) \times \mathrm{G}_1'\) 到 \(\Omega'\)）是 \(\Omega'\) 上某个 \(\mathrm{SL}(n, \mathrm{K})\) Haar 测度的限制。
\item
  \(\Omega'\) 在 \(\mathbf{SL}(n,\mathbf{K})\) 中的补集对于 \(\mathbf{SL}(n,\mathbf{K})\) 的 Haar 测度是零测的。
\end{enumerate}

群 \(\mathbf{GL}(n,\mathbf{K})\) 是幺模的（例 1），而 \(\mathbf{SL}(n,\mathbf{K})\) 是 \(\mathbf{GL}(n,\mathbf{K})\) 的正规子群，因此也是幺模的（ \(\S 2\) ，No.~7, Prop. 10 b)）。断言 b) 由 a)、引理 8 以及 \(\S 2\) ，No.~9 的 Prop. 13 得出。我们来证明 c)。根据引理 8 a)，只需证明如下事实：若 \(p((u_{ij})_{1 \leqslant i,j \leqslant n})\) 是一个在 \(\mathbf{SL}(n,\mathbf{K})\) 上不恒为零的多项式，则满足 \(U \in \mathbf{SL}(n,\mathbf{K})\) 且 \(p(U) = 0\) 的元素所成的集合 E 对 Haar 测度是零测的。考虑 \(\S 1\) ，No.~10, Cor. of Prop. 13，\(\mathbf{SL}(n,\mathbf{K})\) 的拓扑有一个可数基。因此，只需证明：对每个 \(U_0 \in \mathbf{E}\)，存在 \(U_0\) 在 \(\mathbf{SL}(n,\mathbf{K})\) 中的一个邻域，其与 E 的交是零测的；或者说，存在 \(\mathbf{SL}(n,\mathbf{K})\) 中 I 的一个邻域 W，使得 \(U_0^{-1}\mathbf{E} \cap \mathbf{W}\) 是零测的。令 \(\mathbf{W} = \Omega'\)。由 b) 可知，问题归结为证明：对 \((M,N) \in \mathrm{T}_1(n,\mathbf{K}) \times \mathrm{G}_1'\)，满足 \(p(U_0MN) = 0\) 的元素所成的集合对于 \(\mu_1 \otimes \mu_2\) 是零测的。根据 \(\mu_1\) 和 \(\mu_2\) 的表达式（在例 3 和例 5 中计算），这将由以下引理得到：

引理 9. --- 令 \(\psi\) 为 \(\neq 0\) 的 \(\mathbf{K}[\mathbf{X}_1, \ldots, \mathbf{X}_r]\) 中的一个多项式。在空间 \(\mathbf{K}^r\) 中，集合 \(\mathbf{N}\) 由 \(\psi(x_1, \ldots, x_r) = 0\) 定义，对 Haar 测度是零测的。

我们对 r 作归纳论证。当 r = 1 时，引理显然成立，因为此时 N 是有限集。如有必要改变变量的编号，可设 \(\psi \notin K[X_{1}, \ldots, X_{r-1}]\)；写成

\[
\psi \left(\mathrm{X} _ {1}, \dots , \mathrm{X} _ {r}\right) = \mathrm{X} _ {r} ^ {m} \psi_ {0} \left(\mathrm{X} _ {1}, \dots , \mathrm{X} _ {r - 1}\right) + \dots + \psi_ {m} \left(\mathrm{X} _ {1}, \dots , \mathrm{X} _ {r - 1}\right)
\]

其中 \(m > 0\) 且 \(\psi_0 \neq 0\)。在空间 \(\mathbf{K}^{r-1}\) 中，令 \(\mathbf{N}_0\) 为由 \(\psi_0(x_1, \ldots, x_{r-1}) = 0\) 定义的集合，根据归纳假设它是零测的。对于每个 \((x_1, \ldots, x_{r-1}) \notin \mathbf{N}_0\)，\(x_r \in \mathbf{K}\) 中使得 \((x_1, \ldots, x_{r-1}, x_r) \in \mathbf{N}\) 的集合是有限的，因而是零测的。由于 \(\mathbf{K}^r\) 在无穷远处可数（§1, No.~10, Prop. 13 的 Cor.），\(\mathbf{N} \cap [(\mathbf{K}^{r-1} - \mathbf{N}_0) \times \mathbf{K}]\) 在 \(\mathbf{K}^r\) 中是零测的。因此 \(\mathbf{N}\) 是零测的。

例 7. --- Iwasawa 分解 \(\mathbf{GL}(n,\mathbf{K})\).

在这个例子中，K 总表示域 R、C、H 中的一个。若 \(\lambda \in K\)，则规定：当 K = R 时 \(\overline{\lambda}\) 等于 \(\lambda\)，当 K = C 或 H 时取 \(\lambda\) 的共轭。令 E 为 K 上的右向量空间，维数为 n，并令 \(\Phi\) 为 E 上的非退化正厄米型。

引理 10. --- 设 \((f_1, f_2, \ldots, f_n)\) 为 E 的一个基。

\begin{enumerate}
\def\labelenumi{\alph{enumi})}
\item
  存在唯一的 E 的标准正交基 \((e_{1}, e_{2}, \ldots, e_{n})\)，使得 \(f_{i} = e_{1}\alpha_{i1} + e_{2}\alpha_{i2} + \cdots + e_{i}\alpha_{ii}\)（\(i = 1, 2, \ldots, n\)），且对所有 i 都有 \(\alpha_{ii} > 0\)。
\item
  对于固定的 \(\Phi\)，\(e_i\) 和 \(\alpha_{ij}\) 连续依赖于 \((f_1, \ldots, f_n) \in \mathbf{E}^n\)。
\end{enumerate}

令 \(E_{i}=f_{1}K+f_{2}K+\cdots+f_{i}K\)，它的维数为 i。令 \(g_{i}\) 为 \(E_{i}\) 中一个非零元素，正交于 \(E_{i-1}\)，且满足 \(\Phi(g_{i},g_{i})=1\)。对 i 作归纳可见，\((g_{1},\ldots,g_{i})\) 是 \(E_{i}\) 的一个标准正交基。特别地，\((g_{1},\ldots,g_{n})\) 是 E 的一个标准正交基。令 \(\lambda_{i}=\Phi(f_{i},g_{i})\)。由于 \(f_{i}\notin E_{i-1}\)，有 \(\lambda_{i}\neq0\)。置 \(e_{i}=g_{i}\left|\lambda_{i}\right|\lambda_{i}^{-1}\)。则

\[
\Phi (e _ {i}, e _ {i}) = | \lambda_ {i} | ^ {2} \overline {{{{\lambda}}}} _ {i} ^ {- 1} \Phi (g _ {i}, g _ {i}) \lambda_ {i} ^ {- 1} = 1,
\]

因此 \((e_1, \ldots, e_i)\) 也是 \(\mathbf{E}_i\) 的一个标准正交基；此外，\(\Phi(e_i, f_i) = |\lambda_i| \overline{\lambda}_i^{-1} \Phi(g_i, f_i) = |\lambda_i| > 0\)，所以 \(e_i\) 具有 a) 中的性质。令 \((e_1', \ldots, e_n')\) 为具有相同性质的 \(\mathbf{E}\) 的另一个标准正交基。由 \(i\) 作归纳可见，\((e_1', \ldots, e_i')\) 必为 \(\mathbf{E}_i\) 的一个基，因此 \(e_i' = e_i \mu_i\)，其中 \(\mu_i \in \mathbf{K}\)。于是

\[
1 = \Phi (e _ {i} ^ {\prime}, e _ {i} ^ {\prime}) = \overline {{\mu}} _ {i} \Phi (e _ {i}, e _ {i}) \mu_ {i} = \overline {{\mu}} _ {i} \mu_ {i},
\]

且 \(0 < \Phi(e_i', f_i) = \overline{\mu}_i \Phi(e_i, f_i)\)，所以 \(\mu_i > 0\) 且 \(\mu_i^2 = 1\)，从而 \(\mu_i = 1\)，这就证明了 a)。假定已经证明 \(e_i\) 和 \(\alpha_{ij}\) 连续依赖于 \((f_1, \ldots, f_n)\)，其中 \(i < i_0\)，现证明 \(e_{i_0}\) 和 \(\alpha_{i_0j}\) 连续依赖于 \((f_1, \ldots, f_n)\)。当 \(j < i_0\) 时，\(\overline{\alpha}_{i_0j} = \Phi(f_{i_0}, e_j)\) 根据归纳假设连续依赖于 \((f_1, \ldots, f_n)\)。另一方面，

\[
\Phi (f _ {i _ {0}}, f _ {i _ {0}}) = | \alpha_ {i _ {0} 1} | ^ {2} + | \alpha_ {i _ {0} 2} | ^ {2} + \dots + | \alpha_ {i _ {0}, i _ {0} - 1} | ^ {2} + \alpha_ {i _ {0} i _ {0}} ^ {2},
\]

因此 \(\alpha_{i_0i_0}\) 连续依赖于 \((f_1,\ldots ,f_n)\)。因此

\[
e _ {i _ {0}} = (f _ {i _ {0}} - e _ {1} \alpha_ {i _ {0} 1} - \dots - e _ {i _ {0} - 1} \alpha_ {i _ {0}, i _ {0} - 1}) \alpha_ {i _ {0} i _ {0}} ^ {- 1}
\]

连续依赖于 \((f_1, \ldots, f_n)\)。

以下设 \(\mathbf{E} = \mathbf{K}^n\)，并取 \(\Phi\) 为如下形式

\[
\overline {{{x}}} _ {1} y _ {1} + \dots + \overline {{{x}}} _ {n} y _ {n}.
\]

回忆一下，\(\mathbf{U}(n,\mathbf{K})\) 表示相应的酉群。即使 K 不可交换，我们仍将 \(\mathrm{T}_{1}(n,\mathrm{K})\) 记为 \(\mathbf{M}_{n}(\mathrm{K})\) 的上三角矩阵群，其中所有对角元素均为 1。

命题 7. --- 令 \(D_{+}^{*}\) 为对角矩阵的群，其对角元素均 \textgreater0。映射 \((U,D,T)\mapsto UDT\) 是从 \(\mathbf{U}(n,\mathbf{K})\times\mathbf{D}_{+}^{*}\times\mathbf{T}_{1}(n,\mathbf{K})\) 到 \(\mathbf{GL}(n,\mathbf{K})\) 的同胚。

令 \((\varepsilon_1, \ldots, \varepsilon_n)\) 为 \(\mathbf{K}^n\) 的典范基。令 \(X \in \mathbf{GL}(n, \mathbf{K})\)。则 \(X \cdot \varepsilon_i = f_i\) 构成 \(\mathbf{E}\) 的一个基。根据引理 10，可以将这个基 \((f_i)\) 与一个基 \((e_i)\) 联系起来。令 \(U\) 为 \(\mathbf{E}\) 的酉自同构的矩阵，该自同构将 \(\varepsilon_i\) 变为 \(e_i\)。于是

\[
U ^ {- 1} \cdot f _ {i} = \varepsilon_ {1} \alpha_ {i 1} + \varepsilon_ {2} \alpha_ {i 2} + \dots + \varepsilon_ {i} \alpha_ {i i}
\]

其中 \(\alpha_{ii} > 0\) 对 \(i = 1,2,\ldots ,n\) 成立。于是 \(X = UC\)，其中 \(C\) 是矩阵

\[
\left( \begin{array}{c c c c} \alpha_ {1 1} & \alpha_ {2 1} & \ldots & \alpha_ {n 1} \\ 0 & \alpha_ {2 2} & \ldots & \alpha_ {n 2} \\ \vdots & \vdots & \ddots & \vdots \\ 0 & 0 & \ldots & \alpha_ {n n} \end{array} \right).
\]

此外，根据引理 10，U 和 C 连续依赖于 X。另一方面，公式 (8) 表明 C 可写成 DT 的形式，其中 \(D \in D_{+}^{*}\)、\(T \in \mathrm{T}_{1}(n, \mathrm{K})\)，并且 D 和 T 连续依赖于 C。分解 X = UDT 的唯一性来自引理 10 的唯一性。

命题 7 中的同胚称为 \(\mathbf{GL}(n,\mathbf{K})\) 的 Iwasawa 分解。

群 \(\mathbf{G} = \mathbf{D}_{+}^{*} \cdot \mathbf{T}_{1}(n, \mathbf{K})\) 是 K 上对角元素为 \textgreater0 的上三角矩阵的集合。我们将该群的元素 \((z_{ij})\) 与元素

\[
\left(\left(z _ {i i}\right) _ {1 \leqslant i \leqslant n}, \left(z _ {i j}\right) _ {1 \leqslant i <   j \leqslant n}\right) \in \left(\mathbf {R} _ {+} ^ {*}\right) ^ {n} \times \mathrm{K} ^ {n (n - 1) / 2}.
\]

完全仿照例 4 的论证，当 \(\mathbf{K} = \mathbf{R}\) 时，可得该群上的右 Haar 测度为

\[
\left(\prod_ {i = 1} ^ {n} z _ {i i} ^ {- i}\right) \cdot \left(\bigotimes_ {i = 1} ^ {n} d z _ {i i}\right) \otimes \left(\bigotimes_ {i <   j} d z _ {i j}\right).
\]

然后应用 §2, No.~9 的 Prop. 13，可见当把 \(\mathbf{GL}(n,\mathbf{K})\) 与 \(\mathbf{U}(n,\mathbf{K})\times\mathbf{G}\) 通过映射 \((U,S)\mapsto US\) 认同后，\(\mathbf{GL}(n,\mathbf{K})\) 上的一个 Haar 测度为（当 \(\mathbf{K}=\mathbf{R}\) 时）

\[
\left(\prod_ {i = 1} ^ {n} z _ {i i} ^ {- i}\right) \cdot \alpha \otimes \left(\bigotimes_ {i = 1} ^ {n} d z _ {i i}\right) \otimes \left(\bigotimes_ {i <   n} d z _ {i j}\right),\tag{14}
\]

其中 \(\alpha\) 表示 \(\mathbf{U}(n,\mathbf{K})\) 上的一个 Haar 测度。

例 8. --- 厄米型空间。

在这个例子中，K 始终表示域 R、C、H 中的一个。记 \(\delta = \dim_{R} K\)（因而 \(\delta = 1, 2\) 或 4）。右向量空间上的一个厄米型 \(\Phi\)，其空间为 \(K^{n}\)，可写为

\[
\Phi (x, y) = \Phi (x _ {1}, \dots , x _ {n}, y _ {1}, \dots , y _ {n}) = \sum_ {i, j = 1} ^ {n} \overline {{{{x}}}} _ {i} h _ {i j} y _ {j}
\]

其中 \(h_{ij} = \overline{h}_{ji}\) 对所有 i 和 j 都成立。记 H 为 \(\mathbf{M}_{n}(\mathbf{K})\) 的厄米矩阵所成的 R 上向量空间。映射 \((h_{ij}) \mapsto \Phi\) 是 H 到 \(K^{n}\) 上的厄米型向量空间的同构，借此我们认同这两个空间。令 \(H_{+}^{*} \subset H\) 为 \(K^{n}\) 上非退化正厄米型的集合。集合 \(H_{+}^{*}\) 在 H 中是凸的；事实上，若 \(\Phi_{1}, \Phi_{2}\) 属于 \(H_{+}^{*}\)，且 \(\lambda, \mu\) 是满足 \textgreater0 和 \(\lambda + \mu = 1\) 的两个数，则 \(\lambda \Phi_{1} + \mu \Phi_{2}\) 显然是正厄米型；另一方面，若 \((\lambda \Phi_{1} + \mu \Phi_{2})(x, x) = 0\)，则 \(\Phi_{1}(x, x) = \Phi_{2}(x, x) = 0\)，因而 x = 0，所以 \(\lambda \Phi_{1} + \mu \Phi_{2}\) 是非退化的。现在来证明 \(H_{+}^{*}\) 是 H 的开子集。令 S 为满足 \(x = (x_{1}, \ldots, x_{n}) \in K^{n}\) 且 \(x_{1} \overline{x}_{1} + \cdots + x_{n} \overline{x}_{n} = 1\) 的 x 的集合；这是 \(K^{n}\) 的紧子集；若 \(\Phi \in H_{+}^{*}\)，函数 \(x \mapsto \Phi(x, x)\) 在 S 上连续且 \textgreater0，故其下确界 \textgreater0；若 \(\Phi' \in H\) 充分接近 \(\Phi\)，则 \(\Phi'(x, x) > 0\) 对所有 \(x \in S\) 成立，所以 \(\Phi'\) 是正的且非退化。

一般线性群 \(\mathbf{GL}(n,\mathbf{K})\) 在 \(\mathfrak{H}\) 上连续地右作用，即 \((X,\Phi)\mapsto\Phi\circ X\)，也就是按 \((X,H)\mapsto{}^{t}\overline{X}\cdot H\cdot X\) 作用，其中 \(H\) 表示与 \(\Phi\) 对应的厄米矩阵。显然，\(\mathfrak{H}_{+}^{*}\) 在 \(\mathbf{GL}(n,\mathbf{K})\) 下稳定。更确切地，根据 Alg., Ch. IX, §6, No.~1, Th. 1 的 Cor. 1，\(^{1}\) \(\mathfrak{H}_{+}^{*}\) 是 \(\mathbf{GL}(n,\mathbf{K})\) 作用在对应于型 \(\sum_{i=1}^{n}\overline{x}_{i}y_{i}\) 的单位矩阵 \(I_{n}\) 上所得的轨道。这个型的稳定子群是 \(\mathbf{U}(n,\mathbf{K})\)。根据附录 I 的引理 2，可以将 \(\mathfrak{H}_{+}^{*}\) 作为拓扑齐性空间与 \(\mathbf{GL}(n,\mathbf{K})/\mathbf{U}(n,\mathbf{K})\) 认同。

对于每个 \(X \in \mathbf{GL}(n, \mathbf{K})\)，令 \(\widetilde{X}\) 为自同构 \(H \mapsto {}^t \overline{X} \cdot H \cdot X\)，它作用在实向量空间 \(\mathfrak{H}\) 上。若 \(\mu\) 表示加法群 \(\mathfrak{H}\) 的 Haar 测度，则有 \(\widetilde{X}^{-1}(\mu) = |\det \widetilde{X}| \cdot \mu\)（§1, No.~10, Prop. 15 的 Cor. 1）。我们来证明

\[
| \det \widetilde {X} | = | N (X) | ^ {\lambda},\tag{15}
\]

其中 N 表示将 \(\mathbf{M}_{n}(\mathrm{K})\) 视为 R-代数时的范数，且 \(\lambda = 1 - \frac{\delta - 2}{\delta n}\)。只需验证 (15) 对于生成 \(\mathbf{GL}(n,\mathbf{K})\) 的一个生成系成立，因此（A, II, §10, No.~13, Prop. 14 的 Cor. 2）只需考虑以下三类 \(X\)：

\begin{enumerate}
\def\labelenumi{\alph{enumi})}
\tightlist
\item
  \(X\) 是如下形式的映射的矩阵
\end{enumerate}

\[
\left(x _ {1}, \dots , x _ {n}\right) \mapsto \left(x _ {\sigma (1)}, \dots , x _ {\sigma (n)}\right),
\]

其中 \(\sigma \in \mathfrak{S}_n\)。在此情形下，\(X\) 的某个幂等于 1，因此 \(|\det \widetilde{X}| = |\mathrm{N}(X)| = 1\)。

\begin{enumerate}
\def\labelenumi{\alph{enumi})}
\setcounter{enumi}{1}
\tightlist
\item
  \(X\) 是如下形式的映射的矩阵
\end{enumerate}

\[
\left(x _ {1}, \dots , x _ {n}\right) \mapsto \left(a x _ {1}, x _ {2}, \dots , x _ {n}\right).
\]

于是，若 \((h_{ij})\in \mathfrak{H}\)，则 \(\widetilde{X} ((h_{ij})) = (h_{ij}')\)，其中 \(h_{11}' = \overline{a} h_{11}a = |a|^2 h_{11}\)，\(h_{1i}' = \overline{a} h_{1i}\)（\(i > 1\)），并且 \(h_{ij}' = h_{ij}\)（\(i > 1\)，\(j > 1\)）；因此

\[
| \det \widetilde {X} | = | a | ^ {2} | a | ^ {\delta (n - 1)} = | a | ^ {2 + \delta (n - 1)}.
\]

另一方面，若 \(Y = (y_{ij}) \in \mathbf{M}_n(\mathbf{K})\)，则 \(XY = (y_{ij}')\)，其中 \(y_{1j}' = ay_{1j}\)，且 \(y_{ij}' = y_{ij}\)（\(i > 1\)）；因此 \(|\mathrm{N}(X)| = |a|^{\delta n}\)。公式 (15) 再次得到验证。

\begin{enumerate}
\def\labelenumi{\alph{enumi})}
\setcounter{enumi}{2}
\tightlist
\item
  \(X\) 是如下形式的映射的矩阵
\end{enumerate}

\[
\left(x _ {1}, \dots , x _ {n}\right) \mapsto \left(x _ {1} + b x _ {2}, x _ {2}, \dots , x _ {n}\right).
\]

于是 \(\widetilde{X}\big((h_{ij})\big) = (h_{ij}^{\prime})\)，其中 \(h_{11}^{\prime} = h_{11}\)，\(h_{12}^{\prime} = h_{12} + h_{11}b\)，\(h_{1i}^{\prime} = h_{1i}\)（\(i > 2\)），\(h_{22}^{\prime} = h_{22} + \overline{b}h_{12} + \overline{h}_{12}b + \overline{b}h_{11}b\)，\(h_{2i}^{\prime} = h_{2i} + \overline{b}h_{1i}\)（\(i > 2\)），\(h_{ij}^{\prime} = h_{ij}\)（\(i > 2\)，\(j > 2\)）。考虑引理 6，可见 \(|\det \widetilde{X}| = 1\)。同样可验证 \(|\mathrm{N}(X)| = 1\)，这就完成了公式 (15) 的证明。

由此，H 上的测度 \(|\mathrm{N}(H)|^{-\lambda/2}d\mu(H)\) 在 \(\mathbf{GL}(n,\mathbf{K})\) 下不变，因为

\[
\begin{array}{r l} & {\widetilde {X} ^ {- 1} \big (| \mathrm{N} (H) | ^ {- \lambda / 2} d \mu (H) \big) = \big | \mathrm{N} \big (\widetilde {X} (H) \big) \big | ^ {- \lambda / 2} \cdot | \det \widetilde {X} | d \mu (H)} \\ & {\qquad = | \mathrm{N} (H) | ^ {- \lambda / 2} | \mathrm{N} (X) | ^ {- \lambda} | \det \widetilde {X} | d \mu (H) = | \mathrm{N} (H) | ^ {- \lambda / 2} d \mu (H).} \end{array}
\]

若 \(H \in \mathfrak{H}_+^*\)，则 \(H = \widetilde{X}(\mathrm{I}_n) = {}^t\overline{X}X\)，其中 \(X \in \mathbf{GL}(n,\mathbf{K})\)，因此 \(\mathrm{N}(H) = \overline{\mathrm{N}(X)}\mathrm{N}(X) > 0\)。所以在 \(\mathfrak{H}_+^*\) 上，唯一（差一个常数因子）的、在 \(\mathbf{GL}(n,\mathbf{K})\) 下不变的测度（cf.~§2, No.~6, Th. 3）是测度

\[
d \gamma (H) = \mathrm{N} (H) ^ {- \lambda / 2} d \mu (H).
\]

特别地，

\[
d \gamma (H) = (\det H) ^ {- (n + 1) / 2} d \mu (H) \quad \text { when } K = R,
\]

\[
d \gamma (H) = (\det H) ^ {- n} d \mu (H) \quad \text { when } \mathrm{K} = \mathrm{C}.
\]

\appendixheading

*引理 1. --- 令 X 为一个局部紧空间，R 为 X 上的开等价关系，且商空间 X/R 为仿紧；令 \(\pi\) 为从 X 到 X/R 的典范映射。存在 X 上一个连续的实值函数 \(F \geqslant 0\)，使得：

\begin{enumerate}
\def\labelenumi{\alph{enumi})}
\item
  F 关于 R 的每个等价类都不恒为零；
\item
  对于 X/R 的每个紧子集 K，\(\overline{\pi}^{1}(K)\) 与 Supp F 的交是紧的。
\end{enumerate}

对每个点 \(u \in \mathbf{X} / \mathbf{R}\)，选取一个函数 \(f_{u} \in \mathcal{K}_{+}(\mathbf{X})\)，使得 \(f_{u}\) 在 \(\overline{\pi}(u)\) 上不恒为零；令 \(\Omega_{u}\) 为 \(f_{u} > 0\) 的点所成的开集；于是 \(u \in \pi(\Omega_{u})\)。由于 \(\pi\) 是开映射，\(\pi(\Omega_{u})\) 构成 \(\mathbf{X} / \mathbf{R}\) 的一个开覆盖。存在一个局部有限开覆盖 \((\mathrm{U}_{\iota})_{\iota \in \mathbf{I}}\)，它比由 \(\pi(\Omega_{u})\) 构成的覆盖更细，于是（GT, IX, §4, No.~3, Prop. 3）存在单位分解 \((g_{\iota})_{\iota \in \mathbf{I}}\)，它定义在 \(\mathbf{X} / \mathbf{R}\) 上并从属于覆盖 \((\mathrm{U}_{\iota})\)。对每个 \(\iota \in \mathbf{I}\)，选取 \(u_{\iota}\) 使得 \(\mathrm{U}_{\iota} \subset \pi(\Omega_{u_{\iota}})\)。函数 \(\mathrm{F}_{\iota} = (g_{\iota} \circ \pi) \cdot f_{u_{\iota}}\) 属于 \(\mathcal{K}(\mathbf{X})\)，且其支集包含于 \(\overline{\pi}(\mathrm{U}_{\iota})\) 中。因此，\(F_{\iota}\) 的支集构成局部有限族，从而可以定义连续函数 \(F \geqslant 0\)，它在 \(X\) 上：令 \(F = \sum_{\iota \in I} F_{\iota}\)。

对于每个 \(u \in X/R\)，存在某个 \(\iota\) 使得 \(g_{\iota}(u) > 0\)，因而 \(u \in U_{\iota}\)；接着存在 \(x \in \Omega_{u_{\iota}}\) 使得 \(\pi(x) = u\)；于是 \(f_{u_{\iota}}(x) > 0\) 且 \(g_{\iota}(\pi(x)) > 0\)，所以 \(F_{\iota}(x) > 0\)，从而 \(F(x) > 0\)；这证明了 F 具有性质 a)。最后，令 K 为 X/R 的一个紧子集。存在 I 的有限子集 J，使得对于 \(\iota \in I - J\)，有 \(\mathrm{U}_{\iota}\cap \mathrm{K} = \varnothing\)，因而 \(\overline{\pi} (\mathbf{K})\cap \operatorname {Supp}\mathbf{F}_{\iota} = \varnothing\)。于是集合

\[
\overline {{\pi}} ^ {1} (\mathrm{K}) \cap \operatorname{Supp} \mathrm{F} = \overline {{\pi}} ^ {1} (\mathrm{K}) \cap \left(\bigcup_ {\iota \in \mathrm{I}} \operatorname{Supp} \mathrm{F} _ {\iota}\right) = \overline {{\pi}} ^ {1} (\mathrm{K}) \cap \left(\bigcup_ {\iota \in \mathrm{J}} \operatorname{Supp} \mathrm{F} _ {\iota}\right)
\]

是紧的。

引理 2. --- 设 G 为在无穷远处可数的局部紧群，M 为 Baire 空间。假设 G 在 M 上连续且传递地左作用。对每个 \(x \in M\)，令 \(H_{x}\) 为 x 在 G 中的稳定子群，于是映射 \(s \mapsto sx\) 诱导出一个连续双射 \(\varphi_{x}\)，将 \(G/H_{x}\) 映到 M。那么 \(\varphi_{x}\) 是从 \(G/H_{x}\) 到 M 的同胚（换言之（GT, III, §2, No.~5），M 是一个拓扑齐性空间）。

令 \(x_0 \in \mathbf{M}\)。只需证明（同上，Prop. 15），映射 \(s \mapsto sx_0\) 将每个邻域 \(\mathbf{V}\)（其中 \(e\) 在 \(\mathbf{G}\) 中）变为 \(x_0\) 在 \(\mathbf{M}\) 中的一个邻域。令 \(\mathbf{W}\) 为 \(e\) 的对称紧邻域，使得 \(\mathbf{W}^2 \subset \mathbf{V}\)。根据假设，\(\mathbf{G}\) 是一列紧集的并，因而是平移列 \((s_n\mathbf{W})\)（它们是 \(\mathbf{W}\) 的平移）的并。于是 \(\mathbf{M}\) 是一列紧集 \((s_n\mathbf{W}x_0)\) 的并。由于 \(\mathbf{M}\) 是 Baire 空间，存在某个 \(n\)，使 \(s_n\mathbf{W}x_0\) 有一个内点 \(s_nwx_0\)（其中 \(w \in \mathbf{W}\)）。因此 \(x_0\) 是

\[
w ^ {- 1} s _ {n} ^ {- 1} (s _ {n} \mathrm{W} x _ {0}) = w ^ {- 1} \mathrm{W} x _ {0} \subset \mathrm{V} x _ {0},
\]

从而 \(Vx_{0}\) 是 M 中 \(x_{0}\) 的一个邻域。

\appendixheading

*引理 1. --- 令 X、B 为两个局部紧空间，\(\pi\) 为从 X 到 B 的映射，\(\nu\) 为 B 上的正测度。令 \(b \mapsto \lambda_b\)（\(b \in \mathbf{B}\)）为 X 上正测度的一个 \(\nu\)-适当族，使得对每个 \(b \in \mathbf{B}\)，测度 \(\lambda_b\) 集中在 \(\overline{\pi}^1(b)\) 上。置 \(\mu = \int \lambda_b d\nu(b)\)，并假设映射 \(\pi\) 是 \(\mu\)-可测的。

\begin{enumerate}
\def\labelenumi{\alph{enumi})}
\item
  若 \(\mathbf{N} \subset \mathbf{B}\) 是局部 \(\nu\)-零测的，则 \(\overline{\pi}^1(\mathbf{N})\) 是局部 \(\mu\)-零测的。
\item
  若 \(f\) 是 \(\nu\)-可测函数（定义在 \(\mathbf{B}\) 上，值域为某个拓扑空间），则 \(f \circ \pi\) 是 \(\mu\)-可测函数（定义在 \(\mathbf{X}\) 上）。
\end{enumerate}

令 K 为 X 的一个紧子集。我们要证明 \(\overline{\pi}^{1}(N) \cap K\) 是 \(\mu\)-零测的，并且 \(f \circ \pi\) 在 K 上的限制是 \(\mu\)-可测的。现在，K 是一个 \(\mu\)-零测集与一列紧子集 \(K_{n}\) 的并，且 \(\pi|K_{n}\) 连续。只需证明 \(\overline{\pi}^{1}(N) \cap K_{n}\) 是 \(\mu\)-零测的，并且 \(f \circ \pi\) 在 \(K_{n}\) 上的限制是 \(\mu\)-可测的。因此以下不妨假设 \(\pi|K\) 连续。于是 \(\pi(K) = K'\) 是紧的。由于 \(\overline{\pi}^{1}(N) \cap K = \overline{\pi}^{1}(N \cap K') \cap K\)，且 \(N \cap K'\) 是 \(\nu\)-零测的，以下不妨假设 N 是 \(\nu\)-零测的。于是 N 包含于一个 \(\nu\)-零测集 \(N'\) 中，而它是开集的可数交（Ch. IV, §4, No.~6, Th. 4 的 Cor. 2）。由于 \(\pi\) 是 \(\mu\)-可测的，\(\overline{\pi}^{1}(N')\) 是 \(\mu\)-可测的（Ch. IV, §5, No.~5, Prop. 7），所以 \(\overline{\pi}^{1}(N') \cap K\) 是 \(\mu\)-可积的，并且（Ch. V, §3, No.~4, Th. 1）

\[
\mu \big (\overline {{{{\pi}}}} ^ {1} (\mathbf {N} ^ {\prime}) \cap \mathbf {K} \big) = \int_ {\mathbf {B}} \lambda_ {b} \big (\overline {{{{\pi}}}} ^ {1} (\mathbf {N} ^ {\prime}) \cap \mathbf {K} \big) d \nu (b).
\]

现在，若 \(b \notin \mathbf{N}'\)，则集合 \(\overline{\pi}(\mathbf{N}') \cap \mathbf{K}\) 对 \(\lambda_b\) 是零测的，因为 \(\lambda_b\) 集中在 \(\overline{\pi}(b)\) 上。因此 \(\mu\big(\overline{\pi}(\mathbf{N}') \cap \mathbf{K}\big) = 0\)。从而，\(\overline{\pi}(\mathbf{N}) \cap \mathbf{K}\) 是 \(\mu\)-零测的，a) 已得证。另一方面，存在 \(\mathbf{K}'\) 的一个划分，由一个 \(\nu\)-零测集 \(\mathbf{M}\) 和一列紧集 \((\mathbf{K}_n')\) 构成，使得 \(f|_{\mathbf{K}_n'}\) 连续。于是 \(f \circ \pi\) 在每个集合 \(\overline{\pi}^1(\mathbf{K}) \cap \mathbf{K}_n'\) 上的限制是连续的，而 \(\overline{\pi}^1(\mathbf{M}) \cap \mathbf{K}\) 根据 a) 是 \(\mu\)-零测的，因此 \(f \circ \pi\) 在 \(\mathbf{K}\) 上的限制确实是 \(\mu\)-可测的。

\specialchapter{习题}

\exercisesection{}

\begin{enumerate}
\def\labelenumi{\arabic{enumi})}
\item
  设 \(G\) 为紧群。证明每个连续表示 \(\varphi\) 都是从 \(G\) 到 \(\mathbf{R}_+^*\) 的表示，并满足 \(\varphi(G) = \{1\}\)。由此推出 \(G\) 上的相对不变正测度是不变的。
\item
  2. 设 G 为拓扑群，且 G 的换位子群在 G 中稠密。证明 G 到 Hausdorff 交换群的每个连续表示 \(\varphi\) 都满足 \(\varphi(G) = \{e\}\)。由此推出，若 G 局部紧，则 G 上的每个相对不变复测度都是不变的。特别地，G 是幺模的。
\item
  设 \(G\) 为局部紧群，\(\mu\) 为 \(G\) 上的左 Haar 测度，且 \(\nu = \Delta_G^{-1/2} \cdot \mu\)。证明
\end{enumerate}

\[
\boldsymbol {\gamma} (s) \nu = \Delta_ {\mathrm{G}} (s) ^ {1 / 2} \nu , \quad \boldsymbol {\delta} (s) \nu = \Delta_ {\mathrm{G}} (s) ^ {1 / 2} \nu , \quad \stackrel {\vee} {\nu} = \nu .
\]

\begin{enumerate}
\def\labelenumi{\arabic{enumi})}
\setcounter{enumi}{3}
\item
  设 \(G, G'\) 为两个局部紧群，\(V\)（分别为 \(V'\)）为 \(G\)（分别为 \(G'\)）的单位元的一个开邻域，\(\varphi\) 为 \(G'\) 到 \(G\) 的局部同构，定义在 \(V'\) 上，且 \(\varphi(V') = V\)。证明 \(\Delta_G \circ \varphi\) 是 \(\Delta_{G'}\) 在 \(V'\) 上的限制。
\item
  对每个 \(a\)（属于乘法群 \(\mathbf{Q}^*\)），令 \(\varphi(a)\) 为加法群 \(\mathbf{R}\) 的自同构，定义为 \(\varphi(a)x = ax\)。赋予 \(\mathbf{Q}^*\) 离散拓扑。令 \(\mathbf{G}\) 为 \(\mathbf{Q}^*\) 与 \(\mathbf{R}\) 的由 \(\varphi\) 定义的拓扑半直积（GT, III, §2, No.~10）。证明 \(\mathbf{G}\) 局部同构于 \(\mathbf{R}\)，但不是幺模的。
\item
  设 \(\mathbf{G}\) 为含有一个开紧子群 \(\mathbf{H}\) 的局部紧群。证明对每个自同构 \(\varphi\)（它作用于 \(\mathbf{G}\)），有 mod \(\varphi \in \mathbf{Q}_+^*\)。（注意，\(\varphi(\mathbf{H}) \cap \mathbf{H}\) 在 \(\mathbf{H}\) 和 \(\varphi(\mathbf{H})\) 中的指数都是有限的。）证明等式 \(\Delta_{\mathbf{G}}(s) = 1\) 定义了 \(\mathbf{G}\) 的一个包含 \(\mathbf{H}\) 的开子群。
\item
  7. 设 G 为局部紧群，\(\beta\) 为 G 上的左 Haar 测度，A 为 G 的一个子集，B 为 G 的一个相对紧的 \(\beta\)-可积子集，且 \(\beta(\mathrm{B}) > 0\)。证明，若 \(\mathbf{A} \cdot \mathbf{B}\) 的紧子集对于 \(\beta\) 具有有界测度，则 \(\mathbf{A}\) 是相对紧的。（仿照命题 2 的论证。）
\item
  8. 设 G 为局部紧群，A 为 G 的稠密子集，\(\alpha\) 为 G 上的左 Haar 测度，H 为 G 的一个 \(\alpha\)-可测子集，具有如下性质：对每个 \(s \in A\)，\(sH \cap (\mathbb{C}H)\) 和 \(H \cap (\mathbb{C}sH)\) 都是局部 \(\alpha\)-零测的。证明 H 或者是局部 \(\alpha\)-零测的，或者其补集是局部 \(\alpha\)-零测的。（证明 \(\varphi_{H} \cdot \alpha\) 是左不变的。）
\item
  采用第 2 号定理 1 的证明中的记号。令 \(\beta\) 为 G 上满足 \(\beta(f_0) = 1\) 的唯一左 Haar 测度。证明，对每个 \(f \in \mathcal{K}(\mathbf{G})\)，\(\mathrm{I}_g(f)\) 趋于 \(\beta(f)\)，关于 \(\mathfrak{B}\)。（令 \(a\) 为映射 \(g \mapsto \mathrm{I}_g(f)\) 关于 \(\mathfrak{B}\) 的一个聚点。存在一个超滤 \(\mathfrak{U}\)，比 \(\mathfrak{B}\) 更细，使得 \(\mathrm{I}_g(f)\) 趋于 \(a\)，关于 \(\mathfrak{U}\)。另一方面，\(\mathrm{I}_g(f)\) 趋于 \(\beta(f)\)，关于 \(\mathfrak{U}\)。）
\item
  10. 设 G 为局部紧群，\(\mu\) 为 G 上的左 Haar 测度。证明每个 \(\mu\)-可积集 A 都包含于紧集的一个可数并中。（化归到 A 为开集的情形，并注意 G 中存在一个开子群，它是紧子集的可数并。）
\end{enumerate}

¶ 11) 设 G 为在无穷远处可数的非离散局部紧群，β 为 G 上的左 Haar 测度。证明 G 的单位元 e 的每个紧邻域 V 都包含 G 的一个正规子群 H，该子群是 β-零测的，且 G/H 是一个 Polish 局部紧群。（可如下进行：令 L 为由 V 生成的 G 的开子群；令 \(b_{1}, b_{2}, \ldots\) 为关于 L 的左陪集代表元。构造单位元 e 的对称邻域的递减序列 \((\mathrm{V}_{n})\)，使其满足：1) \(\mathrm{V}_{n}^{2} \subset \mathrm{V}_{n-1}\)；2) \(x\mathrm{V}_{n}x^{-1} \subset \mathrm{V}_{n-1}\) 对每个 \(x \in \mathrm{V}\) 成立；3) \(b_{i}\mathrm{V}_{n}b_{i}^{-1} \subset \mathrm{V}_{n-1}\) 对 \(1 \leqslant i \leqslant n\) 成立；4) \(\beta(\mathrm{V}_{n}) \leqslant 1/n\)。置 \(\mathrm{H} = \mathrm{V}_{1} \cap \mathrm{V}_{2} \cap \cdots\)。

令 \((f_{n})\) 为关于 G 的左一致结构一致连续的数值函数序列。证明可以构造 H，使得除此之外，每个 \(f_{n}\) 在 H 的陪集上为常值。（在上述构造中，取 \(V_{n}\) 使得 \(|f_{i}(x) - f_{i}(y)| \leqslant 1/n\) 当 \(x^{-1}y \in V_{n}\) 且 \(1 \leqslant i \leqslant n\) 时。）

令 \((g_n)\) 为 \(\beta\)-可积的数值函数序列，定义在 \(G\) 上。证明可以构造 \(H\)，使得除此之外，每个 \(g_n\) 几乎处处等于一个在 \(H\) 的陪集上为常值的函数。

\begin{enumerate}
\def\labelenumi{\arabic{enumi})}
\setcounter{enumi}{11}
\tightlist
\item
  令 \(\mathbf{K}\) 为非离散局部紧域，\(\mathbf{E}\) 为 \(\mathbf{K}\) 上的左拓扑向量空间。
\end{enumerate}

\begin{enumerate}
\def\labelenumi{\alph{enumi})}
\item
  a) 若 E 是 1 维的，则 E 同构于 \(\mathbf{K}_s\)（仿照 TVS, I, §2, No.~2 的命题 2 的证明，将绝对值替换为函数 mod \(_K\)）。
\item
  b) 若 E 是 n 维的，则 E 同构于 \(K_{s}^{n}\)（仿照 TVS, I, §2, No.~3 的定理 2 的证明）。
\item
  c) 假设 E 局部紧。令 F 为 E 的一个有限维 n 的线性子空间。~对每个 \(a \in K\)，令 \(\text{mod}_{\mathbf{E}}(a)\) 为 E 中映射 \(x \mapsto ax\) 的模。由于根据 b)，F 在 E 中是闭的，可以构造 \(\text{mod}_{\mathbf{E}/\mathbf{F}}(a)\)。证明 \(\text{mod}_{\mathbf{E}}(a) = \text{mod}_{\mathbf{K}}(a)^n \text{mod}_{\mathbf{E}/\mathbf{F}}(a)\)。证明，若 \(\text{mod}_{\mathbf{K}}(a) < 1\) 且 \(F \neq E\)，则 \(\text{mod}_{\mathbf{E}/\mathbf{F}}(a) < 1\)。（有 \(a^n \to 0\)，当 \(n \to +\infty\) 时；由此仿照 No.~10 的命题 12 的论证，推出 \(\text{mod}_{\mathbf{E}/\mathbf{F}}(a^n) \to 0\)。）由此推出 \(n \leqslant \log \text{mod}_{\mathbf{E}}(1/a)/\log \text{mod}_{\mathbf{K}}(1/a)\)，从而 E 是有限维的。
\end{enumerate}

（我们将在 CA, VI 中看到，局部紧域的拓扑可以由一个绝对值定义。于是 c) 的最终结果将成为 TVS, I, §2, No.~4 的定理 3 的特例。）

¶ 13) 设 G 为局部紧群，连续地左作用于局部紧 Polish 空间 T；β 为 G 上的左 Haar 测度，ν 为在 G 下拟不变的 T 上的正测度。

\begin{enumerate}
\def\labelenumi{\alph{enumi})}
\tightlist
\item
  存在函数 \((s,x)\mapsto \chi (s,x) > 0\)，定义在 \(\mathbf{G}\times \mathbf{T}\) 上，它是局部 \((\beta \otimes \nu)\)-可积的，并且对每个 \(s\in \mathbf{G}\)，函数 \(x\mapsto \chi (s^{-1},x)\) 局部 \(\nu\)-可积且满足 \(\gamma(s)\nu = \chi(s^{-1},\cdot)\cdot\nu\)。（证明 \((s,x)\mapsto(s,sx)\) 将 \(\beta\otimes\nu\) 变为等价测度，可利用 Ch. V, §5, No.~5 的定理 2 的准则 3)；令 \(\chi(s^{-1},x)d\beta(s)d\nu(x)\) 为此等价测度。对于 \(\varphi\in\mathcal{H}(\mathrm{G})\) 和 \(\psi\in\mathcal{H}(\mathrm{T})\)，
\end{enumerate}

\[
\int \varphi (s) d \beta (s) \int \psi (s x) d \nu (x) = \int \varphi (s) d \beta (s) \int \psi (x) \chi (s ^ {- 1}, x) d \nu (x),
\]

从而 \(\int \psi(sx) d\nu(x) = \int \psi(x)\chi(s^{-1}, x) d\nu(x)\) 除去一个局部 \(\beta\)-零测集 \(N(\psi)\) 外成立。接着应用 Ch. VI, §3, No.~1 的引理 1；然后在一个零测集上修改 \(\chi\)，该集为 \((\beta \otimes \nu)\)-零测集。

\begin{enumerate}
\def\labelenumi{\alph{enumi})}
\setcounter{enumi}{1}
\tightlist
\item
  证明对于任意的 \(s, t\)（属于 \(\mathbf{G}\)），有
\end{enumerate}

\[
\chi (s t, x) = \chi (s, t x) \chi (t, x)
\]

除去一个 \(\nu\)-零测的 \(x\) 值集合外成立（利用关系 \(\pmb{\gamma}(st)\nu = \pmb{\gamma}(s)(\pmb{\gamma}(t)\nu)\)。c) 证明 a) 中的函数 \(\chi\)（见 \(a\)）由一个局部 \((\beta \otimes \nu)\)-零测子集 \(\mathbf{G} \times \mathbf{T}\) 决定。

¶ 14) 设 T 为局部紧空间，U 为 T 上的一致结构，且存在一个围道 \(V_{0}\)，使得 \(V_{0}(t)\) 对所有 \(t \in T\) 都是紧的（GT, II, §4, Exer. 9）。另一方面，令 \(\Gamma\) 为 T 的同胚群，它关于一致结构 U 一致等度连续，并且存在 \(a \in T\)，其在 \(\Gamma\) 下的轨道稠密。证明 T 上存在一个正测度 \(\neq 0\)，它在 \(\Gamma\) 下不变，并且除一个常数因子外唯一。可按如下方式进行：

\(1^{\circ}\) 关于不变测度的存在性，采用第 2 号定理 1 的方法；先证明：若 K 是 T 的一个紧子集，U 是 \(a\) 的一个开邻域，则存在有限个元素 \(\sigma_{i} \in \Gamma\)，使得 \(\mathrm{K} \subset \bigcup \sigma_{i}(\mathrm{U})\)。

\(2^{\circ}\) 关于唯一性，首先注意，\(\mathcal{U}\) 的围道中在每个 \(\sigma \times \sigma\)（\(\sigma \in \Gamma\)）这个 T × T 的同胚下不变者，构成围道的一个基本系 \(\mathfrak{S}\)。对于 T 的每个相对紧集 A ⊂ T，以及每个具有非空内点的相对紧集 B ⊂ T，令 (A : B) 为用形如 \(\sigma\) B（\(\sigma \in \Gamma\)）的集合覆盖 A 所需的最少集合数；若 C ⊂ T 是第三个具有非空内点的相对紧集，则 (A : C) ≤ (A : B)(B : C)。令 K 为 T 的一个紧子集，令 L ⊃ K 为相对紧开集，令 V ∈ \(\mathfrak{S}\) 为包含于 V0 中的对称开围道，使得 V(K) ⊂ L，并令 W ∈ \(\mathfrak{S}\) 为包含于 V 中的对称闭围道。对每个满足 UW ⊂ V 且 WU ⊂ V 的对称围道 U ∈ \(\mathfrak{S}\)，以及每个 T 上的正测度 \(\nu\)，它在 \(\Gamma\) 下不变，有

\begin{enumerate}
\def\labelenumi{(\arabic{enumi})}
\tightlist
\item
\end{enumerate}

\[
\big (\mathrm{W} (a): \mathrm{U} (a) \big) \nu (\mathrm{K}) \leqslant \big (\mathrm{L}: \mathrm{U} (a) \big) \nu \big (\mathrm{V} (a) \big)\tag{2}
\]

\[
\bigl (\mathrm{K}: \mathrm{U} (a) \bigr) \nu \bigl (\mathrm{W} (a) \bigr) \leqslant \bigl (\mathrm{V} (a): \mathrm{U} (a) \bigr) \nu (\mathrm{L}).
\]

为此，证明：若 \((\sigma_i(\mathrm{U}(a)))\) 是由 \((\mathrm{L}:\mathrm{U}(a))\) 个集合构成的 L 的一个覆盖，则每个 \(x\in \mathbf{K}\) 至少属于 \((\mathrm{W}(a):\mathrm{U}(a))\) 个形如 \(\sigma_{i}(\mathrm{V}(a))\) 的集合，并且每个 \(y\in \mathrm{L}\) 至多属于 \((\mathrm{V}(a):\mathrm{U}(a))\) 个形如 \(\sigma_{i}(\mathrm{W}(a))\) 的集合（注意，若 \(y\in \sigma_i(\mathrm{W}(a))\)，则 \(\sigma_{i}(\mathrm{U}(a))\) 包含于 \(\mathrm{V}(x)\) 中）。

令 \(\mathfrak{F}\) 为一个比 \(\mathfrak{S}\) 的截面滤子更细的超滤；令 \(A_0\) 为 T 中一个非空的相对紧开集，并置 \(\lambda_U(A) = (A: U(a)) / (A_0: U(a))\)，其中 \(U \in \mathfrak{S}\) 为对称围道，A 为 T 的相对紧子集。令 \(\lambda(A) = \lim_{\mathfrak{F}} \lambda_U(A)\)；对 T 的每个紧子集 K，置 \(\lambda'(K) = \inf \lambda(B)\)，其中 B 遍历 K 的相对紧开邻域的集合。由 (1) 和 (2) 推出

\[
\begin{array}{r} \lambda^ {\prime} \big (\mathrm{W} (a) \big) \nu (\mathrm{K}) \leqslant \lambda (\mathrm{L}) \nu \big (\mathrm{W} (a) \big) \\ \lambda^ {\prime} (\mathrm{K}) \nu \big (\mathrm{W} (a) \big) \leqslant \lambda^ {\prime} \big (\mathrm{W} (a) \big) \nu (\mathrm{L}). \end{array}
\]

由此可知，若 \(K_{1}, K_{2}\) 是 T 的两个紧子集，且 \(L_{1} \supset K_{1}\)、\(L_{2} \supset K_{2}\) 是两个相对紧开集，则

\[
\lambda^ {\prime} (\mathrm{K} _ {1}) \nu (\mathrm{K} _ {2}) \leqslant \lambda (\mathrm{L} _ {2}) \nu (\mathrm{L} _ {1}),
\]

最后得到 \(\lambda^{\prime}(\mathrm{K}_{1})\nu (\mathrm{K}_{2}) = \lambda^{\prime}(\mathrm{K}_{2})\nu (\mathrm{K}_{1})\)

¶ 15) 设 G 为局部紧群，H 为 G 的闭子群，从而 G 连续地作用于 G/H。假设 H 满足如下条件：对 G 中 e 的每个邻域 V，存在 e 的一个邻域 U，使得 HU ⊂ VH（注意，如果 e 有一个在 G 的每个内自同构下不变的邻域基本系，则 G 的每个子群 H 都满足此条件）。证明 G/H 上存在一个在 G 下不变的非零正测度。（采用第 2 号定理 1 的证明中的相同方法，并注意：一方面，当 U 遍历 G 中 e 的邻域时，集合 HUH 在典范映射 π : G → G/H 下的像构成 π(H) 的邻域基本系；另一方面，对于 G 中 e 的每个邻域 V 和 G 的每个紧子集 K，存在 G 中 e 的一个邻域 U 和 G 的一个紧子集 L，使得每个与 KH 相交的形如 sHUH 的集合都可写成形如 s'HUH 的集合，其中 s' ∈ L。）

¶ 16) 设 G 为紧交换群，E 为 G 的稠密子集，且在 G 的运算下稳定。假设对 E 上的每个有界数值函数 f，都定义了一个数 M(f)，具有如下性质：对于 E 上任意两个有界函数 f、g，E 中任意元素 \(a_{i}, b_{k}\)（\(1 \leqslant i \leqslant m, 1 \leqslant k \leqslant n\)），以及任意实数 \(\alpha_{i}, \beta_{k}\)（\(1 \leqslant i \leqslant m, 1 \leqslant k \leqslant n\)），若

\[
\sum_ {i} \alpha_ {i} f (x a _ {i}) \leqslant \sum_ {k} \beta_ {k} g (x b _ {k})
\]

对所有 \(x \in \mathbf{E}\) 都有上述不等式，则有不等式

\[
\left(\sum_ {i} \alpha_ {i}\right) \mathrm{M} (f) \leqslant \left(\sum_ {k} \beta_ {k}\right) \mathrm{M} (g).
\]

另假设 \(\mathbf{M}(1) = 1\)（cf.~TVS, IV, Appendix）。

\begin{enumerate}
\def\labelenumi{\alph{enumi})}
\item
  令 \(\mu\) 为 G 上的归一化 Haar 测度。证明，对 G 上的每个连续数值函数 \(f\)，都有 \(\int f d\mu = \mathrm{M}(f|\mathrm{E})\)。（通过取适当的单位分解，先在 E 上考虑恒等于 \(\int f d\mu\) 的常值函数，再用形如 \(x \mapsto \sum_{i} \alpha_i f(xa_i)\) 的函数（其中 \(a_i \in \mathrm{E}\)）去逼近它。）
\item
  b) 由 a) 推出：若对 E 的每个子集 B 置 \(\nu(B) = M(\varphi_B)\)，则对 G 的每个开子集 U，有 \(\nu(U \cap E) \geqslant \mu(U)\)，对 G 的每个闭子集 F，有 \(\nu(F \cap E) \leqslant \mu(F)\)。对于 G 的每个 \(\mu\)-二次集 P（Ch. IV, §5, Exer. 17 d)），有 \(\nu(P \cap E) = \mu(P)\)。
\end{enumerate}

¶ 17) 设 T 为局部紧空间，S 为 T 的一个子集，并赋予其运算 \((x,y)\mapsto xy\)，使 S 成为一个幺半群 \(^{1}\)（不预先假定有单位元）；当 S 取 T 的诱导拓扑时，该运算在 S×S 上连续。再假设 S 的每个元素都是可消去的 \(^{2}\)（A, I, §2, No.~2）。令 \(\mu\) 为 T 上的有界正测度，集中在 S 上（因而 S 是 \(\mu\)-可测的），且总质量为 1。假设 \(\mu\) 在以下意义下左不变：对于 S 上定义的每个连续有界数值函数 f（根据 Ch. IV, §5, No.~5 的 Prop. 8 的 Cor.，因而 \(\mu\)-可测），都有 \(\int f(sx) d\mu(x) = \int f(x)d\mu(x)\)，对所有 \(s \in S\) 成立。这等价于说，对 S 的每个紧子集 K，都有 \(\mu(sK) = \mu(K)\)。

\begin{enumerate}
\def\labelenumi{\alph{enumi})}
\item
  a) 考虑 T×T 上的积测度 \(\mu\otimes\mu\)；证明存在 S 的紧子集 K，使得 K×K 在两个连续映射 \((x,y)\mapsto(x,xy)\) 和 \((x,y)\mapsto(xy,x)\) 下的像的测度任意接近 1（利用 Lebesgue--Fubini 定理）。由此推出这些像有一个公共点，并据此证明 S 有单位元（cf.~A, I, §2, Exer. 9）。
\item
  b) 证明 S 是紧群。（首先证明对每个 \(x \in S\)，内测度 \(\mu_{*}(xS)\) 等于 1，从而 xS 可测且测度集中在 xS 上；xS 随后成为一个可以应用 a) 的结果的幺半群，由此证明 S 是群。仿照命题 2 的论证证明 S 是紧的。最后利用 GT, III, §4.3 的 Exer. 21。）
\end{enumerate}

18) 设 X 为局部紧空间，G 为紧群，连续地作用于 X；E 为 X 中一个点的轨道，\(\mathcal {F}\) 为 X 上连续数值函数所成的向量空间，使得对每个函数 \(f \in \mathcal {F}\) 和每个 \(s \in G\)，都有 \(\gamma (s) f \in \mathcal {F}\)；再假设 \(\mathcal {F}\) 包含 X 上的常值函数。令 \(x _ {0} \in \mathrm{X}\) 为 G 下不变的点，并满足 \(| f (x _ {0}) | \leqslant \sup _ {y \in E} | f (y) |\) 对每个函数 \(f \in \mathcal {F}\) 成立。证明

从而有 \(f(x_0) = \int_{\mathbf{G}} f(s \cdot z) ds\)，对每个 \(z \in \mathbf{E}\) 和每个 \(f \in \mathcal{F}\) 成立。（证明存在一个正测度 \(\nu\)，定义在 \(\mathbf{E}\) 上，总质量为 1，使得 \(f(x_0) = \int_{\mathbf{E}} f(z) d\nu(z)\)，对每个 \(f \in \mathcal{F}\) 成立，然后应用 Lebesgue-Fubini 定理。）*当 \(X = R^n\)、\(G\) 为正交群且 \(x_0 = 0\) 时；应用 §2 的公式 (7)。*

\begin{enumerate}
\def\labelenumi{\arabic{enumi})}
\setcounter{enumi}{18}
\tightlist
\item
  设 X 为紧空间，A 为带单位元的 R 上赋范代数，G 为紧群；假设 G 连续地作用于 A 和 X，且对每个 \(s \in G\)，映射 \(a \mapsto s \cdot a\) 是代数 A 的自同构，并且满足 \(\|s \cdot a\| = \|a\|\) 对所有 \(a \in A\)。称从 X 到 A 的映射 f 在 G 下协变，如果
\end{enumerate}

\[
f (s \cdot x) = s \cdot f (x)
\]

对所有 \(s \in G\) 和 \(x \in X\) 都成立。令 \(B\) 为 \(A^X\) 的一个子环，由协变连续函数组成，并包含所有从 \(X\) 到 \(R\) 的协变连续映射（其中 \(R\) 被认同为 \(A\) 的子代数；注意，对这样的映射 \(g\)，有 \(g(s \cdot x) = g(x)\) 对所有 \(s \in G\) 和 \(x \in X\) 成立）。令 \(f\) 为从 \(X\) 到 \(A\) 的协变连续映射，并假设对每个 \(y \in X\)，存在一个映射 \(g_y \in B\)，使得 \(f(y) = g_y(y)\)。证明，对每个 \(\varepsilon > 0\)，存在一个映射 \(g \in B\)，使得 \(\|f(x) - g(x)\| \leqslant \varepsilon\) 对所有 \(x \in X\) 成立。（利用适当的连续单位分解（\(\varphi_i\)）（定义在 \(X\) 上）并引入函数 \(h_i\)，使 \(h_i(x) = \int_G \varphi_i(s \cdot x) ds\)。）

¶ 20) 设 G 为局部紧群，μ 为 G 上的左 Haar 测度，A 和 B 为 G 的两个子集。

\begin{enumerate}
\def\labelenumi{\alph{enumi})}
\tightlist
\item
  a) 假设下列两个条件之一成立：
\end{enumerate}

\(\alpha)\) A 是 \(\mu\)-可积的；

\(\beta)\ \mu^{*}(A) < +\infty\)，且 \(B\) 是 \(\mu\)-可测的。

证明，在这两种情形中的每一种情形下，函数 \(f(s) = \mu^{*}(sA \cap B)\) 关于 \(G\) 的右一致结构是一致连续的（关于 \(G\)）。(对于 \(M, N\) 为 \(G\) 的任意两个子集，置

\[
d (\mathrm{M}, \mathrm{N}) = \mu^ {*} \big ((\mathrm{M} \cap \mathbb {C N}) \cup (\mathrm{N} \cap \mathbb {C M}) \big).
\]

首先考虑 A 为紧集的情形；利用 Ch. IV, §4, No.~6, Th. 4，证明对每个 \(\varepsilon > 0\)，存在 G 中 \(e\) 的一个邻域 U，使得对每个 \(s \in G\) 和 \(t \in U\)，

\[
d (s \mathrm{A} \cap \mathrm{B}, s t \mathrm{A} \cap \mathrm{B}) \leqslant \varepsilon .
\]

然后应用 Ch. IV, §5 的 Exer. 13。若 B 是 \(\mu\)-可测的，\(\mu^{*}(\mathbf{A}) < +\infty\)，且 \((\mathbf{A}_{n})\) 是 G 的一个包含 A 的可积子集递减序列，满足 \(\inf (\mu (\mathbf{A}_n)) = \mu^{*}(\mathbf{A})\)，证明

\[
\mu^ {*} (s \mathrm{A} \cap \mathrm{B}) = \inf \left(\mu^ {*} (s \mathrm{A} _ {n} \cap \mathrm{B})\right)
\]

（Ch. V, 1st edn., §2, No.~2, Lemma 1）；\(^{4}\) 但注意另一方面，\(\mu(A_{n}-A_{n+1})\) 随 \(1/n\) 趋于 0。）

\begin{enumerate}
\def\labelenumi{\alph{enumi})}
\setcounter{enumi}{1}
\item
  若 \(\mathbf{A}^{-1}\) 是 \(\mu\)-可积的且 \(\mu^{*}(\mathbf{B}) < +\infty\)，则函数 \(f\) 关于 \(\mathbf{G}\) 的右、左一致结构都是一致连续的；此外，有 \(\int_{\mathbf{G}} f(s) d\mu(s) = \mu(\mathbf{A}^{-1})\mu^{*}(\mathbf{B})\)。（化归到 \(\mathbf{B}\) 可积的情形；注意此时 \(\mu(s\mathbf{A} \cap \mathbf{B}) = \mu(\mathbf{A} \cap s^{-1}\mathbf{B})\)，\(\varphi_{s\mathbf{A} \cap \mathbf{B}} = \varphi_{s\mathbf{A}}\varphi_{\mathbf{B}}\)，且 \(\varphi_{s\mathbf{A}}(t) = \varphi_{t\mathbf{A}^{-1}}(s)\)。）
\item
  由 \(a\) 推出，在上述两种情形下，只要 A 和 B 都不是零测的，AB 和 BA 的内部都非空。（Cf. Ch. VIII, §4, No.~6, Prop. 17。）
\item
  d) 在群 \(\mathbf{G} = \mathbf{SL}_{2}(\mathbf{R})\) 中，给出一个紧集 A 和一个 \(\mu\)-可测集 B 的例子，使得 \(f(s) = \mu(sA \cap B)\) 关于左一致结构不是一致连续的。（注意，存在 G 中趋于 e 的元素序列 \((t_{n})\) 和 G 中的元素序列 \((s_{n})\)，使得序列 \((s_{n}^{-1}t_{n}s_{n})\) 趋于无穷远点。）
\item
  e) 在局部紧群 G 中，给出两个不可测集 A、B 的例子，它们的外测度有限，且函数 \(s \mapsto \mu^{*}(sA \cap B)\) 不连续（cf. Ch. IV, §4, Exer. 8）。
\end{enumerate}

\begin{enumerate}
\def\labelenumi{\arabic{enumi})}
\setcounter{enumi}{20}
\tightlist
\item
  \begin{enumerate}
  \def\labelenumii{\alph{enumii})}
  \tightlist
  \item
    设 G 为局部紧群，\(\mu\) 为 G 上的左 Haar 测度。证明，若 S 是 G 的一个稳定子集且 \(\mu_{*}(S)>0\)，则 S 的内部非空（cf.~Exer. 20）。特别地，若 G 的一个子群 H 具有非零内测度，则 H 是 G 的开子群。
  \end{enumerate}
\end{enumerate}

\begin{enumerate}
\def\labelenumi{\alph{enumi})}
\setcounter{enumi}{1}
\item
  若 G 是紧的，则 G 的每个满足 \(\mu_{*}(S) > 0\) 的稳定子集 S 都是 G 的紧开子群（注意，S 的内部是稳定的，并利用 Exer. 17 b)）。
\item
若 G 是紧交换群，采用加法记号，并且 G 中存在一个无限阶元素 a。证明存在 G 的一个稳定子集 S，使 \(a \in S\)、\(-a \notin S\) 且 \(G = S \cup (-S)\)（利用 Zorn 引理）；由 b) 推出 S 不可测且 \(\mu_{*}(S) = 0\)。
\end{enumerate}

\begin{enumerate}
\def\labelenumi{\arabic{enumi})}
\setcounter{enumi}{21}
\tightlist
\item
  设 G 为局部紧群，\(\mu\) 为 G 上的左 Haar 测度，A 为 G 的一个可积子集，且 \(\mu(A) > 0\)。证明集合 H(A) 由满足 \(s \in G\) 且 \(\mu(A) = \mu(A \cap sA)\) 的元素组成，是一个紧群。（利用 Exer. 20，注意 H(A) 是 G 的闭子群。为证明其紧性，取 A 的一个紧子集 B，使 \(\mu(B) > \mu(A)/2\)，并证明 \(H(A) \subset BB^{-1}\)。）
\end{enumerate}

¶ 23) 设 G 为局部紧交换群，采用加法记号；μ 为 G 上的 Haar 测度，A、B 为两个可积子集。

\begin{enumerate}
\def\labelenumi{\alph{enumi})}
\tightlist
\item
  a) 对每个 \(s \in \mathbf{G}\)，令
\end{enumerate}

\[
\mathrm{A} ^ {\prime} = \sigma_ {s} (\mathrm{A}, \mathrm{B}) = \mathrm{A} \cup (\mathrm{B} + s), \quad \mathrm{B} ^ {\prime} = \tau_ {s} (\mathrm{A}, \mathrm{B}) = (\mathrm{A} - s) \cap \mathrm{B}.
\]

证明 \(\mu (\mathrm{A}^{\prime}) + \mu (\mathrm{B}^{\prime}) = \mu (\mathrm{A}) + \mu (\mathrm{B})\) 且 \(\mathrm{A}' + \mathrm{B}'\subset \mathrm{A} + \mathrm{B}\)

\begin{enumerate}
\def\labelenumi{\alph{enumi})}
\setcounter{enumi}{1}
\tightlist
\item
  假设 0 属于 \(\mathbf{A} \cap \mathbf{B}\)。称 \((\mathbf{A}', \mathbf{B}')\)（\(\mathbf{G}\) 的可积子集）由 (A, B) 导出，如果存在元素的序列 \((s_k)_{1 \leqslant k \leqslant n}\)，其元素属于 \(\mathbf{G}\)，和两个子集序列 \((\mathbf{A}_k)_{0 \leqslant k \leqslant n}\)、\((\mathbf{B}_k)_{0 \leqslant k \leqslant n}\)，它们属于 \(\mathbf{G}\)，使得 \(\mathbf{A}_0 = \mathbf{A}\)、\(\mathbf{B}_0 = \mathbf{B}\)、\(\mathbf{A}_k = \sigma_{s_k}(\mathbf{A}_{k-1}, \mathbf{B}_{k-1})\)、\(\mathbf{B}_k = \tau_{s_k}(\mathbf{A}_{k-1}, \mathbf{B}_{k-1})\) 对 \(1 \leqslant k \leqslant n\) 成立，\(s_k \in \mathbf{A}_{k-1}\) 对 \(1 \leqslant k \leqslant n\) 成立，且 \(\mathbf{A}' = \mathbf{A}_n\)、\(\mathbf{B}' = \mathbf{B}_n\)。证明存在一列有序对 \((\mathrm{E}_n, \mathrm{F}_n)\)，使得 \(\mathbf{E}_0 = \mathbf{A}\)、\(\mathbf{F}_0 = \mathbf{B}\)、\((\mathrm{E}_{n+1}, \mathrm{F}_{n+1})\) 由 \((\mathrm{E}_n, \mathrm{F}_n)\) 导出，并且 \(\mu((\mathrm{E}_n - s) \cap \mathrm{F}_n) \geqslant \mu(\mathrm{F}_{n+1}) - 2^{-n}\) 对每个 \(n\) 和每个 \(s \in \mathrm{E}_n\) 成立。置 \(\mathbf{E}_{\infty} = \bigcup_{n} \mathbf{E}_n\)、\(\mathbf{F}_{\infty} = \bigcap_{n} \mathbf{F}_n\)。证明对每个 \(s \in \mathrm{E}_{\infty}\)，
\end{enumerate}

\[
\mu \big ((\mathrm{E} _ {\infty} - s) \cap \mathrm{F} _ {\infty} \big) = \mu (\mathrm{F} _ {\infty}).
\]

\begin{enumerate}
\def\labelenumi{\alph{enumi})}
\setcounter{enumi}{2}
\tightlist
\item
  假设 \(\mu (\mathrm{F}_{\infty}) > 0\)。证明函数
\end{enumerate}

\[
f (s) = \mu \big ((\mathrm{E} _ {\infty} - s) \cap \mathrm{F} _ {\infty} \big)
\]

它只能取 0 和 \(\mu(\mathrm{F}_{\infty})\) 这两个值；并且集合 C 由满足 \(s \in G\) 且 \(f(s) = \mu(\mathrm{F}_{\infty})\) 的点组成，既开又闭，满足 \(\mu(C) = \mu(\mathrm{E}_{\infty})\)，且是 \(\mathrm{E}_{\infty}\) 的闭包（利用 Exer. 20 a) 和 b)）。另一方面，令 D 为满足 \(s \in \mathrm{F}_{\infty}\) 且 \(\mathrm{F}_{\infty}\) 与 \(s\) 的每个邻域的交的测度为 \(>0\) 的点所成的集合。证明 \(\mu(D) = \mu(\mathrm{F}_{\infty})\)，并且

\[
\mathbf {E} _ {\infty} + \mathbf {D} \subset \mathbf {C};
\]

由此推出 D 包含于 Exer. 22 所定义的子群 H(C)，且 H(C) 是 G 的紧开子群。最后证明 C + H(C) = C，\(\mu(C) \geqslant \mu(A) + \mu(B) - \mu(H(C))\)，以及 C ⊂ A + B（对每个 c ∈ C，考虑 E∞ ∩ (c - F∞) 的测度）。

\begin{enumerate}
\def\labelenumi{\alph{enumi})}
\setcounter{enumi}{3}
\tightlist
\item
  由 c) 推出，对于 G 的两个可积子集 A、B：或者 \(\mu_{*}(A+B)\geqslant\mu(A)+\mu(B)\)；或者存在 G 的一个紧开子群 H，使得 \(A+B\) 包含 H 的一个陪集，此时 \(\mu_{*}(A+B)\geqslant\mu(A)+\mu(B)-\mu(H)\)。考虑 G 连通的情形。
\end{enumerate}

¶ 24 a) 在 R 中，令 A（分别为 B）为具有如下性质的数 x 的集合：它们的正规二进展开 \(x = x_0 + \sum_{i=1}^{\infty} x_i 2^{-i}\)（\(x_0\) 为整数，\(x_i = 0\) 或 \(x_i = 1\)，当 \(i \geqslant 1\) 时）满足 \(x_{i}=0\) 对 i 为偶数且 \textgreater0（分别地，i 为奇数时满足此条件）。证明，对于 Lebesgue 测度，A 和 B 的测度为零，但 \(A+B=R\)。

\begin{enumerate}
\def\labelenumi{\alph{enumi})}
\setcounter{enumi}{1}
\item
  由 a) 推出，R 中存在一个包含于 A ∪ B 的基 H（在 Q 上），因而其测度为零。数字 rh（其中 r ∈ Q、h ∈ H）组成的集合 P₁ 也具有测度零。
\item
  令 \(P_{n}\) 为这样的实数集合：相对于基 H，其坐标中至多有 n 个非零坐标。证明，若 \(P_{n}\) 是零测的且 \(P_{n+1}\) 可测，则 \(P_{n+1}\) 是零测的。（令 \(h_{0} \in H\)；先证明，\(x \in P_{n+1}\) 中相对于 \(h_{0}\) 的坐标为 \(\neq 0\) 的那些点所成的集合 S 是零测的。利用 Exer. 20 证明，若 \(P_{n+1}\) 不是零测的，则存在两个不同的点 \(x', x''\)，它们属于 \(P_{n+1} \cap CS\)，使得 \((x' - x'') / h_{0}\) 是有理数，并由此导出矛盾。）
\item
  d) 由 a) 和 b) 推出，R 中存在两个零测集 C、D，使得 C + D 不可测。
\end{enumerate}

¶ 25) a) 令 f 为定义在 R 上的正数值函数，它关于 R 上的 Lebesgue 测度 μ 可积、有界且支集紧。令 γ = sup f(t)。

对每个 \(w \in \mathbf{R}\)，记 \(\mathrm{U}_f(w)\) 为 \(t \in \mathbf{R}\) 中满足 \(f(t) \geqslant w\) 的点所成的集合；置 \(\nu_f(w) = \mu^*(\mathrm{U}_f(w))\)。证明，对每个 \(\alpha > 1\)，

\[
\int_ {- \infty} ^ {+ \infty} f ^ {\alpha} (t) d t = \int_ {0} ^ {\gamma} \nu_ {f} (w) \alpha w ^ {\alpha - 1} d w.
\]

\begin{enumerate}
\def\labelenumi{\alph{enumi})}
\setcounter{enumi}{1}
\tightlist
\item
  令 g 为满足与 f 相同条件的另一个数值函数，并置 \(\delta = \sup_{t \in \mathbf{R}} g(t)\)。令 h 为定义在 \(R^{2}\) 上的函数：\(h(u,v) = f(u) + g(v)\) 当 \(f(u)g(v) \neq 0\) 时，否则 \(h(u,v) = 0\)；最后置 \(k(t) = \sup_{u+v=t} h(u,v)\)，于是 k 是正的、可积的、有界的且支集紧。证明，对每个 \(\alpha > 1\)，
\end{enumerate}

\[
\int_ {- \infty} ^ {+ \infty} k ^ {\alpha} (t) d t \geqslant (\gamma + \delta) ^ {\alpha} \left(\frac {1}{\gamma^ {\alpha}} \int_ {- \infty} ^ {+ \infty} f ^ {\alpha} (t) d t + \frac {1}{\delta^ {\alpha}} \int_ {- \infty} ^ {+ \infty} g ^ {\alpha} (t) d t\right).
\]

（注意，当 \(0 < w < 1\) 时，有 \(\mathrm{U}_k(\gamma w + \delta w) \supset \mathrm{U}_f(\gamma w) + \mathrm{U}_g(\delta w)\)，并利用 \(a\) 和 Exer. 23 d)。）

\begin{enumerate}
\def\labelenumi{\alph{enumi})}
\setcounter{enumi}{2}
\tightlist
\item
  c) 令 \(\mu_{n}\) 为 \(\mathbf{R}^{n}\) 上的 Lebesgue 测度，A、B 为 \(\mathbf{R}^{n}\) 的两个可积子集。证明（Brunn-Minkowski 不等式）
\end{enumerate}

\[
\left(\left(\mu_ {n}\right) _ {*} (\mathrm{A} + \mathrm{B})\right) ^ {1 / n} \geqslant \left(\mu_ {n} (\mathrm{A})\right) ^ {1 / n} + \left(\mu_ {n} (\mathrm{B})\right) ^ {1 / n}.
\]

（化归到 A 和 B 都是紧集的情形。然后利用 Exer. 23 d)、Lebesgue--Fubini 定理、b) 中证明的不等式以及 Hölder 不等式，对 n 作归纳论证。）

¶ 26) 设 G 为局部紧群，μ 为 G 上的左 Haar 测度，k 为满足 ≥1 的整数，A 为可积集。证明，对每个 ε \textgreater{} 0，存在 G 中 e 的一个邻域 U，具有如下性质：对于 U 中含有 k 个元素的每个有限子集 S，满足 sS ⊂ A 的 s ∈ G 所成的集合的测度至少为 (1 - ε)μ(A)。（化归到 A 为紧集的情形，并取 U 使得

\[
\mu (\mathrm{AU}) \leqslant \left(1 - \frac {\varepsilon}{k - 1}\right) \mu (\mathrm{A}).
\]

对 G 的每个有限子集 H，置 \(p(\mathrm{H}) = \operatorname{Card}(\mathrm{H})\)，并置 \(q(\mathrm{H}) = 1\)（当 \(\mathrm{H} \neq \emptyset\) 时），以及 \(q(\mathrm{H}) = 0\)（当 \(\mathrm{H} = \emptyset\) 时）；通过计算积分

\[
\int_ {\mathbf {G}} p (\mathrm{A} \cap s \mathrm{S}) d s \quad \text { and } \quad \int_ {\mathbf {G}} q (\mathrm{A} \cap s \mathrm{S}) d s  ,
\]

证明：若 \(h \leqslant k\)，令 \(\mathbf{M}_h\) 为 \(s \in G\) 中满足 \(\operatorname{Card}(\mathbf{A} \cap s\mathbf{S}) = h\) 的元素所成的集合，则

\[
\sum_ {h = 1} ^ {k} h \mu (\mathrm{M} _ {h}) = k \cdot \mu (\mathrm{A}) \quad \text { and } \quad \sum_ {h = 1} ^ {k} \mu (\mathrm{M} _ {h}) \leqslant \mu (\mathrm{AU}).
\]

由此推出 \(\sum_{h=1}^{k-1} h\mu(M_h) \leqslant k\varepsilon \cdot \mu(A)\)。

\begin{enumerate}
\def\labelenumi{\arabic{enumi})}
\setcounter{enumi}{26}
\tightlist
\item
  设 \(\mathbf{G}\) 为作用于集合 \(\mathbf{X}\) 上的群。称 \(\mathbf{P}\)（分别为 \(\mathbf{C}\)）为 \(\mathbf{X}\) 的子集，并将其称为 G-填充集（分别为 G-覆盖集），如果对每个 \(s \neq e\)（\(\mathbf{G}\) 中的），都有 \(s\mathbf{P} \cap \mathbf{P} = \emptyset\)（分别地，如果 \(\mathbf{X} = \bigcup_{s \in \mathbf{G}} s\mathbf{C}\)）。称同时为 G-填充集和 G-覆盖集的子集 \(\mathbf{P}\) 为 G-铺砌集。
\end{enumerate}

\begin{enumerate}
\def\labelenumi{\alph{enumi})}
\item
  假设 X 局部紧，G 可数，连续地作用于 X，并且 X 上存在一个在 G 下不变的非零正测度 \(\mu\)。令 P、C 为 \(\mu\)-可积的 G-填充集和 G-覆盖集。证明 \(\mu(C) \geqslant \mu(P)\)。（注意，\(\mu(C) \geqslant \sum_{s \in G} \mu(C \cap sP)\)。）
\item
  假设此外 X 上存在一个与 X 的拓扑相容的一致结构，并且有一个由 G 不变的开围道组成的基本系 S。另一方面，令 \(\Delta(G)\) 为 X 的所有可积 G-覆盖集 C 的 \(\mu(C)\) 的下确界。令 V 为属于 S 的一个围道，令 \(a \in X\) 满足 \(\mu(V(a)) > \Delta(G)\)；证明存在 \(s \in G\)，使得 \(s \neq e\) 且 \((a, sa) \in \overline{V} \circ V\)。
\item
  假设 X 是局部紧群，\(\mu\) 是左 Haar 测度，G 是 X 的一个可数子群，并通过左平移作用。沿用对 \(\Delta(G)\) 的定义（见 \(b\)），证明，若 A 是 X 的可积子集且 \(\mu(A) > \Delta(G)\)，则存在 \(s \in G \cap AA^{-1}\)，使 \(s \neq e\)。特殊情形：若 \(X = R^n\)，且 G 是秩为 \(n\) 的离散子群，则 \(\mathbf{R}^n: \Delta(G)\) 等于 G 关于 Z 的一个基相对于 \(\mathbf{R}^n\) 的典范基的行列式绝对值；若 A 是内部非空的对称闭凸集 \(^5\)（在 \(\mathbf{R}^n\) 中），且 \(\mu(A) \geqslant 2^n \Delta(G)\)，证明 \(A \cap G\) 中存在一个异于 0 的点（Minkowski 定理）。
\item
在 a) 的假设下，令 \(f\) 为一个 \(\geqslant 0\) 且 \(\mu\)-可积的定义在 \(X\) 上的函数。证明存在 \(a, b\) 属于 \(X\) 的两个点，使得 \(\mu(C) \sum_{s \in G} f(sa) \geqslant \int_X f(x) d\mu(x)\) 且 \(\mu (\mathrm{P})\sum_{s\in \mathbb{G}}f(sb)\leqslant \int_{\mathbb{X}}f(x)d\mu (x).\)（注意，若 \(g\) 是一个可积函数 \(\geqslant 0\)，E 是
\end{enumerate}

\(\mu\)-可积集 \(\mathbf{X}\) 中，则存在 \(c \in \mathbf{E}\)，使得 \(\int_{\mathbf{E}} g(x) d\mu(x) \leqslant g(c)\mu(\mathbf{E})\)，并存在 \(c' \in \mathbf{E}\)，使得 \(\int_{\mathbf{E}} g(x) d\mu(x) \geqslant g(c')\mu(\mathbf{E})\)。

¶ 28) 在 Exer. 27 a) 的假设下，再假设存在一个 μ-可积的 G-铺砌集 F。

\begin{enumerate}
\def\labelenumi{\alph{enumi})}
\tightlist
\item
对于定义在 X 上的每个 \(\mu\)-可积数值函数 \(f \geqslant 0\)，置 \(\widetilde{f}(x) = \sum_{s \in G} f(sx)\)，从而 \(\int_{F} \widetilde{f}(x) d\mu(x) = \int_{X} f(x) d\mu(x)\)（§2, No.~10, Prop. 15）。令 \((f_i)_{1 \leqslant i \leqslant n}\) 为数值函数 \(\geqslant 0\)、\(\mu\)-可积于 X 的一个族；当 \(i \neq j\) 时，置
\end{enumerate}

\[
m _ {i j} = \int_ {\mathbf {F}} \widetilde {f} _ {i} (x) \widetilde {f} _ {j} (x) d \mu (x),
\]

并令 \(c_{i} = \sup_{x\in X}\widetilde{f}_{i}(x)\)。证明至少存在一对不同的指标 \((i,j)\)，使得

\[
m _ {i j} \geqslant \frac {1}{n (n - 1) \mu (\mathrm{F})} \left(\left(\sum_ {i = 1} ^ {n} \int_ {\mathrm{X}} f _ {i} (x) d \mu (x)\right) ^ {2} - \mu (\mathrm{F}) \sum_ {i = 1} ^ {n} c _ {i} \int_ {\mathrm{X}} f _ {i} (x) d x\right).
\]

（利用 Cauchy-Schwarz 不等式从下方估计和 \(\sum_{i\neq j}m_{ij}\)。）由此推出：若 \((\mathrm{A}_i)_{1\leqslant i\leqslant n}\) 是 \(\mu\)-可积的 G-填充集的有限族，且

\[
\sum_ {i = 1} ^ {n} \mu (\mathrm{A} _ {i}) > \mu (\mathrm{F}),
\]

则存在一对不同的指标 \((i,j)\) 以及 \(s\in G\)，使得 \(\mu (\mathrm{A}_is\cap \mathrm{A}_j) > 0\)。b) 若 \(\mathbf{G}_0\) 是 \(G\) 的有限指数子群，且 \((G:\mathbf{G}_0) = h\)，且 \((s_1,\ldots ,s_h)\) 是 \(\mathbf{G}_0\) 在 \(G\) 中的右陪集代表元系，则 \(\mathrm{F_0} = \bigcup_{1\leqslant i\leqslant h}s_i\mathrm{F}\) 是一个 \(G_{0}\)-铺砌集。

\begin{enumerate}
\def\labelenumi{\alph{enumi})}
\setcounter{enumi}{2}
\tightlist
\item
  在 \(\mathbf{X} = \mathbf{R}^n\) 中，令 \(\mathbf{A}\) 为内部非空的对称闭凸集；令 \(u_i : (x_j) \mapsto \sum_{j=1}^{n} c_{ij} x_j\) 为 \(m\) 个定义在 \(\mathbf{X}\) 上的线性型，其整数系数为 \(c_{ij} (m < n)\)。
\end{enumerate}

证明，对每个整数 \(p \geqslant 1\) 以及每个数 \(r \geqslant 0\)，若满足 \(\mu(A)r^n \geqslant 2^n p^m\)，则存在一个异于 0 的点 \(x \in rA \cap Z^n\)，并且 \(u_i(x) \equiv 0 \pmod{p}\) 对 \(1 \leqslant i \leqslant m\) 成立（将 Minkowski 定理（Exer. 27 c)）应用于子群 \(G_0\)，它是 \(Z^n\) 的子群，由 \(z \in Z^n\) 中满足 \(u_i(z) \equiv 0 \pmod{p}\) 对 \(1 \leqslant i \leqslant m\) 成立的元素构成，并利用 \(b\)）。特殊情形：当 \(n = 2\)、\(m = 1\) 且 A 由 \(|x_1| \leqslant 1\)、\(|x_2| \leqslant 1\) 定义时（Thue 定理）。

¶ 29) a) 令 p 为素数；存在两个整数 a、b，使 \(a^{2}+b^{2}+1 \equiv 0 \pmod{p}\)（Alg., Ch. V, 1st edn., §11, No.~5, Cor. of Th. 3）。证明存在不全为零的整数 \(x_{1}, x_{2}, x_{3}, x_{4}\)，使 \(ax_{1} + bx_{2} \equiv x_{3} \pmod{p}\)、\(bx_{1} - ax_{2} \equiv x_{4} \pmod{p}\)，并且

\[
y = x _ {1} ^ {2} + x _ {2} ^ {2} + x _ {3} ^ {2} + x _ {4} ^ {2} \leqslant \sqrt {2} p
\]

（采用 Exer. 28 c) 中的相同方法。）证明 \(y\) 能被 \(p\) 整除，并由此推出 \(y = p\)。

\begin{enumerate}
\def\labelenumi{\alph{enumi})}
\setcounter{enumi}{1}
\tightlist
\item
  b) 由 a) 推出，每个整数 \(n \geqslant 0\) 都是四个平方数之和（Lagrange 定理；利用 Alg., Ch. IV, 1st edn., §2, Exer. 11）。
\end{enumerate}

\begin{enumerate}
\def\labelenumi{\arabic{enumi})}
\setcounter{enumi}{29}
\tightlist
\item
  设 G 为局部紧群，\(\mu\) 为 G 上的左 Haar 测度，\(\nu\) 为 G 上的有界测度。假设映射 \(s \mapsto \gamma(s)\nu\) 从 G 到 Banach 空间 \(\mathcal{M}^{1}(G)\) 连续。证明测度 \(\nu\) 以 \(\mu\) 为基。（令 A 为 G 的一个 \(\mu\)-零测紧子集。仿照命题 11 的论证，证明对稠密子集中的 x 有 \(\nu(xA)=0\)。由此推出 \(\nu(A)=0\)。）
\end{enumerate}

反过来，若 \(\nu\) 以 \(\mu\) 为基，则映射 \(s \mapsto \pmb{\gamma}(s)\nu\) 从 G 到 \(\mathcal{M}^1(\mathrm{G})\) 连续：cf.~Ch. VIII, §2, No.~5, Prop. 8。
